% Options for packages loaded elsewhere
\PassOptionsToPackage{unicode}{hyperref}
\PassOptionsToPackage{hyphens}{url}
\PassOptionsToPackage{dvipsnames,svgnames,x11names}{xcolor}
%
\documentclass[
  11pt,
]{article}
\usepackage{amsmath,amssymb}
\usepackage{setspace}
\usepackage{iftex}
\ifPDFTeX
  \usepackage[T1]{fontenc}
  \usepackage[utf8]{inputenc}
  \usepackage{textcomp} % provide euro and other symbols
\else % if luatex or xetex
  \usepackage{unicode-math} % this also loads fontspec
  \defaultfontfeatures{Scale=MatchLowercase}
  \defaultfontfeatures[\rmfamily]{Ligatures=TeX,Scale=1}
\fi
\usepackage{lmodern}
\ifPDFTeX\else
  % xetex/luatex font selection
  \setmainfont[BoldFont=lmroman10-bold.otf,ItalicFont=lmroman10-italic.otf,BoldItalicFont=lmroman10-bolditalic.otf]{lmroman10-regular.otf}
  \setsansfont[BoldFont=lmsans10-bold.otf,ItalicFont=lmsans10-oblique.otf,BoldItalicFont=lmsans10-boldoblique.otf]{lmsans10-regular.otf}
  \setmonofont[ItalicFont=lmmono10-italic.otf,HyphenChar=None]{lmmono10-regular.otf}
  \setmathfont[]{latinmodern-math.otf}
\fi
% Use upquote if available, for straight quotes in verbatim environments
\IfFileExists{upquote.sty}{\usepackage{upquote}}{}
\IfFileExists{microtype.sty}{% use microtype if available
  \usepackage[]{microtype}
  \UseMicrotypeSet[protrusion]{basicmath} % disable protrusion for tt fonts
}{}
\makeatletter
\@ifundefined{KOMAClassName}{% if non-KOMA class
  \IfFileExists{parskip.sty}{%
    \usepackage{parskip}
  }{% else
    \setlength{\parindent}{0pt}
    \setlength{\parskip}{6pt plus 2pt minus 1pt}}
}{% if KOMA class
  \KOMAoptions{parskip=half}}
\makeatother
\usepackage{xcolor}
\usepackage[margin=1in]{geometry}
\usepackage{longtable,booktabs,array}
\usepackage{calc} % for calculating minipage widths
% Correct order of tables after \paragraph or \subparagraph
\usepackage{etoolbox}
\makeatletter
\patchcmd\longtable{\par}{\if@noskipsec\mbox{}\fi\par}{}{}
\makeatother
% Allow footnotes in longtable head/foot
\IfFileExists{footnotehyper.sty}{\usepackage{footnotehyper}}{\usepackage{footnote}}
\makesavenoteenv{longtable}
\setlength{\emergencystretch}{3em} % prevent overfull lines
\providecommand{\tightlist}{%
  \setlength{\itemsep}{0pt}\setlength{\parskip}{0pt}}
\setcounter{secnumdepth}{-\maxdimen} % remove section numbering
\newlength{\cslhangindent}
\setlength{\cslhangindent}{1.5em}
\newlength{\csllabelwidth}
\setlength{\csllabelwidth}{3em}
\newlength{\cslentryspacingunit} % times entry-spacing
\setlength{\cslentryspacingunit}{\parskip}
\newenvironment{CSLReferences}[2] % #1 hanging-ident, #2 entry spacing
 {% don't indent paragraphs
  \setlength{\parindent}{0pt}
  % turn on hanging indent if param 1 is 1
  \ifodd #1
  \let\oldpar\par
  \def\par{\hangindent=\cslhangindent\oldpar}
  \fi
  % set entry spacing
  \setlength{\parskip}{#2\cslentryspacingunit}
 }%
 {}
\usepackage{calc}
\newcommand{\CSLBlock}[1]{#1\hfill\break}
\newcommand{\CSLLeftMargin}[1]{\parbox[t]{\csllabelwidth}{#1}}
\newcommand{\CSLRightInline}[1]{\parbox[t]{\linewidth - \csllabelwidth}{#1}\break}
\newcommand{\CSLIndent}[1]{\hspace{\cslhangindent}#1}
\usepackage{longtable,booktabs,array}
\usepackage{calc}
\usepackage{needspace}
\usepackage{placeins}
\usepackage{xurl}
\usepackage{seqsplit}
\emergencystretch=3em
\brokenpenalty=10000
\widowpenalty=10000
\clubpenalty=10000
% Pages are ragged-bottom, so this stretch is never visible; it lets a page end between tight-list items.
\renewcommand{\tightlist}{\setlength{\itemsep}{0pt plus 1pt}\setlength{\parskip}{0pt}}
\setlength{\tabcolsep}{4pt}
\renewcommand{\arraystretch}{1.12}
\setlength{\parindent}{0pt}
\setlength{\parskip}{5pt plus 1pt minus 1pt}
\renewcommand{\topfraction}{0.9}
\renewcommand{\textfraction}{0.08}
\renewcommand{\floatpagefraction}{0.8}
\setlength{\textfloatsep}{14pt plus 4pt minus 2pt}
\AtBeginDocument{\hypersetup{breaklinks=true}}
\ifLuaTeX
  \usepackage{selnolig}  % disable illegal ligatures
\fi
\IfFileExists{bookmark.sty}{\usepackage{bookmark}}{\usepackage{hyperref}}
\IfFileExists{xurl.sty}{\usepackage{xurl}}{} % add URL line breaks if available
\urlstyle{same}
\hypersetup{
  pdftitle={Supplement: signals, interfaces and worked-example data},
  pdfauthor={Jonathan Matkowsky},
  colorlinks=true,
  linkcolor={black},
  filecolor={Maroon},
  citecolor={Blue},
  urlcolor={blue},
  pdfcreator={LaTeX via pandoc}}

\author{}
\date{}

\begin{document}

\setstretch{1}
\FloatBarrier

\begingroup\raggedright\hyphenpenalty=10000\exhyphenpenalty=10000

\hypertarget{supplement-signals-interfaces-and-worked-example-data}{%
\section{Supplement: signals, interfaces and worked-example
data}\label{supplement-signals-interfaces-and-worked-example-data}}

\endgroup

\suppressfloats[t]

Jonathan Matkowsky

10 November 2026 · Accepted working draft for version 1.1 · Not
published · Not peer reviewed

This supplement accompanies the \href{MANUSCRIPT.md}{complete
manuscript}. S1 and S2 explain the sources and selection. S3 to S5
specify the signals, records and interfaces. S6 supplies the
hypothetical investigations and their arithmetic; S7 and S8 give the
proposed evaluations. S9 extends the comparison with prior work. The 19
CSV files use \mbox{UTF-8} encoding and stable mechanism, source and
signal identifiers. Citation labels match the package's
\href{REFERENCES.bib}{BibLaTeX references}.

\FloatBarrier

\hypertarget{s1.-sources-and-comparisons}{%
\subsection{S1. Sources and
comparisons}\label{s1.-sources-and-comparisons}}

\suppressfloats[t]

\href{SUPPLEMENT/S1_CITED_SOURCE_NOTES.csv}{Source notes} identify the
version, supporting passages, publication status and supported claim for
each of the 122 cited works. The
\href{SUPPLEMENT/S1_REQUIREMENTS_MATRIX.csv}{requirements matrix}
compares 50 sources against C1 to C3. It codes two parts of C1 and C2
and three of C3: an outcome linked to its decision, and late,
unobservable and false reports kept apart (C1); query coverage and
learned time (C2); typed separation of observations, inferences and
memory, a human principal before reuse, and links from records an agent
output informed back to that output (C3). E means explicit support, P
partial or adjacent support, and a dot means the cited material does not
establish the requirement. The \texttt{payment\_setting} column records
yes for a payment or payment-fraud setting, P for a setting partly in
payments or adjacent to payment fraud, and no otherwise. It marks two
cyber threat intelligence sources, STIX 2.1 and Hettema's
belief-revision paper, as cyber; other security and threat-intelligence
sources are coded no. The source notes provide passage locators, linked
to the matrix by citation key. Their inspection dates are in UTC.

\href{SUPPLEMENT/S1_STAGE_COVERAGE.csv}{Stage coverage} supplies the
source-level codes behind manuscript Table 1. Each family receives the
strongest code among its sources, so a family row may combine
capabilities from several systems. Source notes record bibliographic and
access issues separately from claims about a method.

Study populations and denominators constrain the comparisons. Sharma's
balanced synthetic sample contains 60 cases drawn from transactions its
classifier scored between 0.2 and 0.8: thresholding the classifier at
0.5 decides 43 correctly and the agent 39. Six of the agent's eight
overrides replace a correct decision with an error. Fatima's
verification rates concern genuine participants under the study's
completion and exclusion rules; they do not measure fraud detection.
Doerfler's 47.8 percent same-session success rate concerns challenged
sign-ins. Its 97.3 percent eventual-access rate counts users, not
sign-ins: it is the share of users with access within one week.
Doerfler's Section 5.3 and Table 7 caption take users who failed a
challenge as the denominator; Doerfler's introduction takes users who
saw a challenge. Neither Fatima's nor Doerfler's aggregate rate
establishes whether an individual event was decided correctly.

FTC complaints describe allegations, and the stipulated orders do not
establish those allegations on their merits. Vendor publications support
attributed descriptions of interfaces or observations, not independent
evidence of industry-wide effectiveness.

\FloatBarrier

\hypertarget{s2.-product-scope-and-mechanism-selection}{%
\subsection{S2. Product scope and mechanism
selection}\label{s2.-product-scope-and-mechanism-selection}}

\suppressfloats[t]

\href{SUPPLEMENT/S2_SURFACE_BASELINE.csv}{Product scope} contains 88
public product and feature entries, read on 27 September 2026: 53
catalog items, twenty items from the Q3 2026 public roadmap (seventeen
appear in a cached copy of that roadmap; the other three are confirmed
as shipped in Stripe's changelog), ten separately documented
capabilities and five careers-page footer items. These entries define
the scope for candidate selection, not a complete or current product
catalog. Each row lists candidate mechanism IDs, not validated coverage
or retained signal mappings; a dot means none is listed. The entries
list 85 of the 99 candidates; the other 14 appear only in the allocation
table. The allocation table gives each candidate's final disposition,
and S3 supplies the retained mappings. Six retained mechanisms are not
listed against any product entry: \mbox{ID-14} (account farming and
aged-account preparation for later abuse), \mbox{PAY-04} (card testing
through attacker-controlled merchant or connected accounts),
\mbox{MP-02} (payment-method cashing through a self-controlled merchant
or connected account), \mbox{MP-04} (merchant misrepresentation and
transaction laundering behind a front), \mbox{MP-07} (merchant-side
dispute evasion) and \mbox{BOT-04} (CAPTCHA and challenge defeat). S5
identifies the exact API schemas used for interface comparisons.

A retained mechanism targets a documented capability, serves a plausible
adversary goal and has supporting public evidence. Sources include
research, incident reports, government and intergovernmental guidance,
and platform documentation. Four public job postings for abuse, risk and
machine-learning roles helped prioritize topics but do not establish
incidents; the manuscript cites two (Section 3), and the allocation keys
record the other two. The stated priority basis names a posting topic
for 24 candidates. Twenty-three of them are among the 26 candidates
allocated Tier 1; all 23 are mapped and carry 65 of the 163 links. The
24th, \mbox{EC-06}, is unmapped. Every mapped mechanism has support
beyond postings in the mechanism-to-source links.

\href{SUPPLEMENT/S2_MECHANISM_ALLOCATION.csv}{Mechanism allocation}
records all 99 candidates and their final dispositions. The signal map
retains 61 mechanisms: 52 of the 53 dedicated candidates, all eight
targeted candidates and the one adjacent-logged candidate, \mbox{ID-18},
insider access abuse by platform or partner staff. Of the other 38,
thirteen were judged covered by existing records, four were screened
without identified supporting evidence and logged as low-priority gaps,
twenty low-priority candidates were not reviewed, and one reviewed
candidate, \mbox{PAY-12}, has no signal mapping. It concerns
triangulation fraud: filling real customers' orders through another
seller using stolen payment credentials. Its review record notes
framework evidence only.
\href{SUPPLEMENT/S2_ALLOCATION_KEYS.csv}{Allocation keys} decode the
source abbreviations. The \texttt{Class} code, also the prefix before
each key's colon, gives the source type: a for law-enforcement
announcements and vendor or incident reports, including EchoLeak
(a:ECHO), which is cited by its arXiv identifier; b for research papers;
b* for other arXiv preprints; c for government and intergovernmental
publications and the MITRE fraud framework; d for the illustrative
platform's documentation, roadmap, blog and public repositories; and e
for the illustrative platform's job postings (Stripe careers). All four
b keys name a peer-reviewed venue. In the \texttt{Verification} column,
V marks a key recorded as verified, and S marks a key that rests on a
summarizing retrieval whose exact passage was not retained, as the
\texttt{d:S} row states. The allocation file repeats S as a suffix on
its evidence keys, as in \mbox{\texttt{a:SSN\ (S)}}. Priority labels
explain selection emphasis; they do not measure evidence quality.

\hypertarget{reading-the-selection-records}{%
\paragraph{Reading the selection
records}\label{reading-the-selection-records}}

The allocation file's \texttt{disposition} field states the final
selection decision. Signal coverage is recorded separately in S3; a
reviewed candidate need not have a signal mapping.

\begingroup\small

\begin{longtable}[]{@{}
  >{\raggedright\arraybackslash}p{(\columnwidth - 2\tabcolsep) * \real{0.2000}}
  >{\raggedright\arraybackslash}p{(\columnwidth - 2\tabcolsep) * \real{0.8000}}@{}}
\toprule\noalign{}
\begin{minipage}[b]{\linewidth}\raggedright
Disposition
\end{minipage} & \begin{minipage}[b]{\linewidth}\raggedright
Meaning
\end{minipage} \\
\midrule\noalign{}
\endhead
\bottomrule\noalign{}
\endlastfoot
dedicated & Reviewed as a distinct mechanism \\
targeted & Received a focused review of its implications for the
design \\
covered & Judged covered by existing records rather than retained as a
separate mapped mechanism \\
screened gap & Screened without identified supporting evidence and
logged as a low-priority gap \\
unreviewed & Not reviewed; all twenty such candidates have low
priority \\
adjacent logged & Recorded as adjacent to the core mechanisms; the one
such candidate is retained in the signal map \\
\end{longtable}

\endgroup

\mbox{\texttt{Priority\ (basis)}} and \mbox{\texttt{Impact\ (basis)}}
retain the selection judgments and their stated reasons. In a priority
basis, \texttt{r:} introduces a topic named in a job posting, with the
posting's key in parentheses where given, and \texttt{p:} introduces
platform-documentation keys. A posting key such as \texttt{(EM)} is the
allocation key \texttt{e:EM} without its class prefix;
platform-documentation keys keep their \texttt{d:} prefix. Impact bases
use no prefixes. \texttt{Scope} distinguishes core and adjacent
mechanisms. \texttt{Parent} records a stated relationship to another
mechanism; a dot means none is listed. \mbox{\texttt{Evidence\ keys}}
resolve to source descriptions in the allocation-keys file. They are
identifiers, not evidence grades. \texttt{Maturity} is the estimate,
before review, of how developed the published research was;
\texttt{after\_review\_maturity} pairs it with the later classification
where both were supplied. \texttt{Evidence} is the confidence label used
to allocate effort (verified, provisional or unresolved), not a
certification of each source's claims. \texttt{Depth} is the initial
allocation: Tier 1 for the deeper review, Tier 2 otherwise, or a logged
gap when evidence was unresolved. \texttt{recorded\_status} states each
review's outcome. The product-scope file links documented capabilities
to candidates; it does not supply a product-availability assessment.

Searches covered nine abuse paths, coded U1 to U9: API-key takeover
(U1), identity provenance at onboarding (U2), document integrity,
including identity and business documents (U3), account compromise (U4),
payout redirection (U5), card testing (U6), abusive merchants (U7), mule
and pass-through funds (U8), and attacks on agents or operator
assistants (U9). Only the entity-relationship notes cite these codes.
Searches followed relevant publications, public guidance and references.
This focused selection does not establish exhaustive coverage of abuse
or related systems.

\hypertarget{s2.1-claim-types}{%
\subsubsection{S2.1 Claim types}\label{s2.1-claim-types}}

Research can propose an untested design, while documentation can
directly establish an interface. Neither publication status nor the
number of sources validates the integrated specification. No claim in
either document rests on author experience.

\begin{table}[htp]
\textbf{Table S2.1. How to read the manuscript's claims.} Claim type
describes what a statement establishes. Source type and publication
status describe where its support comes from.\par\smallskip
{\small
\begin{tabular}{@{}
  >{\raggedright\arraybackslash}p{(\columnwidth - 2\tabcolsep) * \real{0.2500}}
  >{\raggedright\arraybackslash}p{(\columnwidth - 2\tabcolsep) * \real{0.7500}}@{}}
\toprule\noalign{}
\begin{minipage}[b]{\linewidth}\raggedright
Claim type
\end{minipage} & \begin{minipage}[b]{\linewidth}\raggedright
What it means
\end{minipage} \\
\midrule\noalign{}

Fact & A description supported by evidence, including a public
specification or an artifact this study produced; source and version are
given where they matter \\
Peer-reviewed prior research & A finding reported in a peer-reviewed
publication whose publication status was confirmed, applicable only to
the setting and conditions it studied \\
Inference & A conclusion motivated by stated premises; the sources do
not establish the conclusion directly \\
Design requirement & A property chosen for the proposed platform on the
stated evidence and reasoning; it is not claimed as necessary for every
platform or as observed in a deployment \\
Hypothesis & A proposition to test for the proposed platform, including
every detection candidate in the manuscript \\
Proposed implementation & One possible way to satisfy a requirement \\
Empirically validated result & A measurement, valid for the proposition
and setting actually tested; results for individual components do not
validate this integrated design \\

\bottomrule\noalign{}
\end{tabular}
}
\end{table}

\hypertarget{s2.2-external-coverage-check}{%
\subsubsection{S2.2 External coverage
check}\label{s2.2-external-coverage-check}}

An external check compares the candidates with MITRE's Fight Fraud
Framework (F3), version 1.1
{[}\protect\hyperlink{ref-MITRE2026F3}{1}{]}. The
\href{SUPPLEMENT/S2_F3_CROSSWALK.csv}{F3 crosswalk} maps each of its 123
techniques and sub-techniques to the candidate mechanisms that cover it,
with a reason for each scope decision. Of the 92 within the product
scope, a retained mechanism covers 34 fully and 40 in part. Of the other
18, fourteen are covered only by candidates outside the signal map. Five
of the fourteen rest only on candidates judged covered by existing
records, among them device fingerprint spoofing; six rest only on
unreviewed candidates, among them 3-D Secure bypass; and three rest on
candidates of both kinds. None rests on a screened gap. Four are covered
by no candidate: DLL injection in the browser, abuse of accessibility
features, theft from password managers and SEO poisoning. Scope and
coverage are one reviewer's judgments, recorded for each row.

\FloatBarrier

\hypertarget{s3.-signals-mappings-and-lineage}{%
\subsection{S3. Signals, mappings and
lineage}\label{s3.-signals-mappings-and-lineage}}

\suppressfloats[t]

The \href{SUPPLEMENT/S3_SIGNALS.csv}{signal inventory} defines all 71
signals, with their families, classes, feature inputs, mechanism and
source IDs, public counterparts, entities, uses and limits. Its
\texttt{internal\_representation\_status} value, DR for every signal,
marks a design requirement rather than observed platform
instrumentation.
\href{SUPPLEMENT/S3_MECHANISM_SIGNAL_PAIRS.csv}{Mechanism-to-signal
mappings} contain 163 links: 154 design mappings motivated by evidence
and nine marked as inferential. These links explain why a record is
proposed; they do not establish a detector's effectiveness or make its
fields uniquely necessary.

The outcome-channel requirements use \mbox{SIG-70} and \mbox{SIG-71} to
record each payment-level or account-level control decision's effect on
each outcome channel, \mbox{SIG-27} to separate authentication results
from reported liability status, and \mbox{SIG-58} for dispute reasons
and payout returns. \mbox{SIG-70} also keeps manual-review dispositions.
Each record retains its source, version and uncertainty.

The inventory contains 26 source observations (SO), two normalized
observations (NO), 24 derived features (DF), two model or risk outputs
(MO) and ten control outcomes (CO). Seven signals provide external
knowledge (XK). The inventory therefore covers five of the seven record
classes in manuscript Section 5.1, plus external knowledge. Analytical
inferences and finished intelligence are later outputs, not input
signals.

\mbox{SIG-44}, \mbox{SIG-45} and \mbox{SIG-46} apply across mechanisms.
They record technique-catalog mappings, minimum observation provenance
and the provenance of test records, respectively. The other 68 signals
map to specific mechanisms. All 61 mapped mechanisms have at least one
mapping motivated by evidence; S2 explains \mbox{PAY-12}'s exclusion.

Each input ID for the 24 derived features names another signal in the
inventory. No feature depends, directly or indirectly, on itself. For
example, \mbox{SIG-24} uses \mbox{SIG-27} to enumerate its eligible
transaction cohort, including completed payments with no report. The
separate chargeback-or-alert-resolution measure also uses \mbox{SIG-25}.
Keeping these inputs explicit makes the denominator and the distinction
between measures inspectable.

The \href{SUPPLEMENT/S3_SOURCE_DIRECTORY.csv}{source directory}
identifies 61 design sources. Of these, 27 are not among the 122 cited
works and have no S1 source note. Each source was read in full in
October 2026. For one, a vendor blog post (L112), the cited passages
were read verbatim and the other paragraphs only as a close paraphrase.
One FTC complaint (L127) was first read in part on 21 October; a
complete copy was read on 29 October. The directory records each
source's type, the depth and version of the read, with qualifiers such
as the version read or pages not read, the access date and the scope.
Where only a preprint or an author copy was read, the row says so.
\href{SUPPLEMENT/S3_MAPPING_SOURCE_JOINS.csv}{Mechanism/source joins}
supply 126 links, each with its source identifier, recorded depth and
scope limit. These records preserve the premises behind the mappings,
including limits on transferring evidence from another setting.

The public-counterpart classification asks whether a pinned schema
contains an object, field or event that supplies part of the required
observation. Of the 71 signals, 32 have partial counterparts, the
question does not apply to 32, and none was identified for seven. S5
gives the exact locators and explains the limits of partial
counterparts; the inventory records each signal's classification.

Identifier links establish limited relationships. A
storefront-to-merchant identifier link does not show that a payment
through the merchant account settled. A verified identifier from the
same namespace and relevant interval can support a stronger, or
\emph{hard}, link. A descriptor or shared IP address may support only a
weaker, or \emph{soft}, link. Neither alone establishes common control
or malicious intent.

The inventory preserves the manuscript's interpretation rules:
onboarding-link histories are a design analogy (Section 5.2), and error
estimates retain their cohorts and denominators (Section 7.2). Each
outcome channel retains its decision context (Section 7.1), and each
broader outcome measure requires its own evaluation (Sections 8.2 and
9.1).

\textbf{Exposure and timing measurements (manuscript Section 5.2).} Meli
and colleagues found a median of about 1,800 unique candidate secrets
daily. Of the monitored secrets, 81 percent remained sixteen days after
discovery. They did not test credentials against the services
{[}\protect\hyperlink{ref-Meli2019}{2}{]}. Demir and colleagues
confirmed 1,748 active credentials for 14 services, with payment
services among the largest groups. Traced through nine archived crawl
snapshots, the verified credentials had been publicly exposed for 12
months on average {[}\protect\hyperlink{ref-Demir2026}{3}{]}. That
duration does not show continuous credential validity. In a deliberate
cloud-key exposure experiment, an actor used the key within minutes
while an automatic quarantine limited its actions. Researchers later
removed that restriction to observe further operations
{[}\protect\hyperlink{ref-Gamazo2023}{4}{]}. Microsoft reported
attackers launching follow-on payment fraud attempts as little as five
minutes after stealing credentials and session cookies through phishing
that relayed the victim's login
{[}\protect\hyperlink{ref-Microsoft2022AiTM}{5}{]}.

\FloatBarrier

\hypertarget{s4.-responsibilities-entities-and-interfaces}{%
\subsection{S4. Responsibilities, entities and
interfaces}\label{s4.-responsibilities-entities-and-interfaces}}

\suppressfloats[t]

The \href{SUPPLEMENT/S4_FUNCTIONS.csv}{functional responsibilities} file
defines 24 responsibilities without requiring separate processes or
databases. Seventeen were derived with payment-specific mechanisms or
sources in view. One (23) adapts agent-security research, and six (2, 3,
4, 13, 21 and 22) come from cyber threat intelligence work. The file's
\texttt{class} column marks these three groups DI, AGV and INH.

Responsibilities 16 and 17 cover engineering validation and adaptive
testing of detections, controls and models. Responsibility 21
independently reviews all agent outputs, including intelligence drafts,
hunting and detection candidates, and recommendations; responsibility 22
provides human release. Permission to investigate grants no authority
over accounts, payments or controls. An authorized human may decide
urgent containment before quality review (manuscript Section 6.4),
without releasing the draft into organizational memory.

The \href{SUPPLEMENT/S4_ENTITIES.csv}{entity inventory} and
\href{SUPPLEMENT/S4_ENTITY_RELATIONS.csv}{entity relationships} define
the entities and the kinds of links between them. Each link records its
source, identifier namespaces and the basis for associating each
identifier with its entity. It retains its effective interval, learned
time, version, confidence, confounders (other explanations for the
link), and permitted purpose and scope. A shared attribute alone does
not establish that one person controls both entities. When a decision
makes an outcome unobservable, the record describes the channel's
status; it does not invent an outcome report.
\protect\hyperlink{complete-storage-requirements}{S4.1} gives the
complete storage requirements.

\Needspace{150pt}

The query interface must show what each query actually covered and what
its result means:

\begingroup\small

\begin{longtable}[]{@{}
  >{\raggedright\arraybackslash}p{(\columnwidth - 2\tabcolsep) * \real{0.2500}}
  >{\raggedright\arraybackslash}p{(\columnwidth - 2\tabcolsep) * \real{0.7500}}@{}}
\toprule\noalign{}
\begin{minipage}[b]{\linewidth}\raggedright
Part of the interface
\end{minipage} & \begin{minipage}[b]{\linewidth}\raggedright
Information or behavior required
\end{minipage} \\
\midrule\noalign{}
\endhead
\bottomrule\noalign{}
\endlastfoot
Request & Task and question; authorized requester; tenant, purpose and
data scope; entity namespaces; event interval and as-of cutoff (only
data the system had learned by then); schema and query versions; row,
time and cost limits; expected observation and required coverage \\
Response & Query/run ID; scope actually evaluated; source and schema
versions; event and ingestion watermarks (how far event time and data
loading had progressed); known gaps and retention horizon for each
source; sampling and access filtering; whether the result is complete,
partial, truncated, denied or unavailable, and any scope-limited flag;
meaning of any count; returned rows or approved aggregate, with the
identifiers of the returned records kept for the run; lineage and
coverage limits \\
Interpretation & A positive observation does not establish how many
cases exist in the population. An empty result counts against a
hypothesis only when the abuse should have produced the queried
observation and the records cover where it would appear. A denied or
unavailable result does not establish adversary behavior. \\
Failure & Enforce limits outside the agent's language model. Stop or
narrow the query after it exceeds a row, time or cost limit, or after an
authorization error. Record incomplete tests and any abstention (a
decision not to conclude). Do not treat missing or not-assessed data as
low risk. \\
\end{longtable}

\endgroup

Some permitted operations also need approval bound to the exact request
payload, target, task, tenant, tool and tool-manifest version. The
manifest is the declared list of the agent's tools. Each approval has an
expiry and a rule governing reuse. At execution, the enforcement gateway
(manuscript Section 6.5) checks that the operation matches the approval
and remains permitted. The restricted agent-and-approval ledger records
both approval and execution.

Control and outcome records distinguish the decision, its inputs and
version, what happened next, the observable channels, actual reports and
any review of correctness. Population error estimates retain their own
cohort, denominator, outcome window, uncertainty and scope; they do not
establish correctness for individual events. Late reports, unobservable
outcomes and false reports have separate fields because they can
coexist. Report presence, channel state and query access also remain
separate. Each temporary delay retains its decision and effective
interval. A derived fact's learned time is when it became queryable,
kept separately from its inputs' latest availability.

A hypothetical timeline illustrates the as-of cutoff. Freeze a
prediction on day 0 using only information then known. A day-30 report
may evaluate it within a stated outcome window, but a model trained on
day 20 cannot use that report. A later training run may use it if the
run's cutoff and label rules permit. Likewise, a relationship learned on
day 10 cannot enter day-0 features merely because its effective date is
earlier.

The \href{SUPPLEMENT/S4_DRAFT_SCHEMA.csv}{draft schema} supplies the
fields that the rule checker, quality reviewers and the releaser use to
assess an intelligence draft. S7 explains the evaluation conditions;
S7.1 gives hypotheses and criteria for E1 to E9. S8 gives the
practitioner review protocol.

\textbf{Returned records a draft does not cite (manuscript Section
6.4).} The checker identifies every record that the case's query runs
returned within a claim's scope but that the draft neither cites nor
lists against it. It attaches those that a fixed rule marks as able to
count against the claim, such as reports, reversals and records of the
benign explanations the claim must exclude, and summarizes the rest by
type and count. Reviewers and the releaser see both beside the claim.
Full checks of claims about named entities examine every marked record
and a recorded random sample of the rest, whose size is fixed in advance
for each claim type, and E8 reports how often that sample contained a
record that counted against the claim. If a share \(q\) of the unmarked
records count against the claim, a random sample of \(m\) finds at least
one with probability at least \(1-(1-q)^m\): 59 records give 95 percent
when \(q\) is 5 percent.

\hypertarget{complete-storage-requirements}{%
\subsubsection{S4.1 Complete storage
requirements}\label{complete-storage-requirements}}

Manuscript Table 3 summarizes this complete storage contract; manuscript
Section 5.2 gives the motivating evidence. Each block groups
requirements by the information it preserves. The stores distinguish
record meanings and need not be separate databases.

\needspace{4\baselineskip}

\textbf{Observation log}

\begin{itemize}
\tightlist
\item
  Preserve records and revisions without overwriting earlier records.
  Keep source, event and delivery identifiers, schema version, entities,
  source lineage, and actor and request identifiers where available. A
  record delivered more than once takes effect only once.
\item
  Record each observation's environment (production, test, synthetic,
  replay or emulation), marked when the record is created, and its
  tenant scope: the business or other boundary within which it may be
  accessed and used. Keep each record's source trust mark. Mark content
  that businesses, customers or outside sources supplied as untrusted.
  Both marks survive normalization and storage.
\item
  For content fetched from a source the subject controls, such as a
  storefront or business website, record the vantage point of each
  fetch: network origin, region and client. A subject can serve
  different content to different visitors, including reviewers.
\item
  Mark every label, assessment, control outcome, draft, proposed lesson
  and released item that a test, replay or emulation creates, or that
  derives from such a record or from synthetic data, with the same
  environment value when it is created. Keep these records out of
  production labels, intelligence and lessons. Record any authorization
  for their evaluation use, with who gave it, the evaluation and its
  scope.
\item
  Keep event or effective time (when something happened or took effect)
  with its source and uncertainty bound. Record the source system's
  recorded time and the platform's ingestion time separately.
\item
  For each store, retain when each record and revision became visible to
  queries.
\item
  A derived record keeps its inputs' latest learned time separately from
  the time it first became queryable. Its own learned time is the latter
  and cannot precede any input's learned time.
\end{itemize}

\needspace{4\baselineskip}

\textbf{Entity-relation graph}

\begin{itemize}
\tightlist
\item
  Label each identifier link by kind and as hard (stronger) or soft
  (weaker). Retain its source, identifier namespaces, effective
  interval, learned time, version, confidence and possible benign
  explanations (confounders).
\item
  Record why each identifier is believed to belong to its entity.
\item
  Keep versions of linked-entity clusters and their membership rules.
  Record end customers separately for each business. Record each link's
  permitted purpose and scope. A link across businesses also records its
  authorization for that identifier type (manuscript Sections 5.2 and
  5.5).
\end{itemize}

\needspace{4\baselineskip}

\textbf{Feature store}

\begin{itemize}
\tightlist
\item
  Keep inputs traceable and values recomputable, with the feature
  definition version, the entities and the computation window.
  Decision-time features use only information available at that
  decision.
\item
  A feature computed over a cluster or other threat-knowledge item
  records that item's version, writer and status. If an agent proposed
  the item, or the item links to an agent output or released
  organizational-memory item that informed it, the feature links to that
  output or item and counts as dependent on it.
\item
  Join later outcomes as evaluation labels under an explicit outcome
  window and maturity period. The outcome window is the period after the
  decision in which an outcome must occur to count.
\end{itemize}

\needspace{4\baselineskip}

\textbf{Label store}

\begin{itemize}
\tightlist
\item
  \textbf{Report history:} retain report type, source, attributed
  transaction or entity and lag. Keep three times: when the reported
  outcome occurred, where known, with its uncertainty; when the report
  was made at its source; and when it arrived at the platform. Preserve
  each value and its learned time, such as a dispute being filed,
  answered, then won or lost, and record whether each won or lost status
  followed adjudication. Append each status change, reversal or
  correction as a new record that references the earlier one. Never edit
  an earlier record, so that what the platform believed at each time can
  be reconstructed.
\item
  \textbf{Truth assessment:} link any assessment that the report is
  mistaken or deliberately false. Keep that assessment in the assessment
  store with its model or reviewer identity and version, how the report
  was selected for review and, for a sampled review, its selection
  probability.
\item
  \textbf{Dependencies:} link every decision that may have shaped the
  report, and any agent output or released organizational-memory item
  that informed the label. An outcome channel is a route, such as a
  chargeback, through which an outcome can be reported. A report may
  link to several decisions, such as a blocked attempt, a retry and the
  approval of the paid attempt, and one decision, such as an account
  restriction, may link to many reports. Each link states the decision's
  role, such as a decision on the reported item, on an earlier related
  attempt, or on an account or payout that contains it, and the basis
  for the link. These links do not change how a measure counts reports
  (Graph semantics). Record the writer's case or session and the agent
  outputs and released items open in it when the record was written.
\item
  \textbf{Channel status:} for each affected transaction, payout or
  account application, keep report presence, channel state and query
  access separate at each read time. These statuses may be computed when
  read from the stored decision, report and channel records. When a
  draft cites a status, store that status with an identifier and its
  read time.
\item
  \textbf{State history:} a channel is open, temporarily delayed or not
  yet opened, permanently closed by a recorded decision, reopened or
  extended by a recorded rule event, or unknown. Keep every transition's
  effective interval and the decision ID, or the rule event's source and
  version, that caused it.
\item
  \textbf{No report:} an open channel remains pending before its
  reporting window ends. It matures with no report only if it stayed
  open throughout the window and the records cover that window. If a
  recorded decision or rule event changed the channel's meaning, the
  matured absence carries that decision's ID, or the rule event's source
  and version, and is read under the changed meaning, not as a negative
  for the original outcome. A channel with no defined window never
  supports a negative by maturity alone. A matured absence is
  provisional: it records the rule version under which it matured and
  returns to pending if a later rule event reopens the window.
\item
  \textbf{Unavailable records:} mark inaccessible records unavailable
  without overwriting known channel state or earlier reports.
\item
  \textbf{Clocks and maturity:} for a temporary delay, keep the expected
  and actual clock-start events when known. Store the reporting window
  and clock-start event with the channel, including their source and
  version. Each channel record also names its rail, the attributes its
  window depends on, such as whether the debited account is a consumer
  or business account or the type of claim, and whether the channel
  adjudicates report truth. Store the maturity used by an evaluation
  with that evaluation.
\end{itemize}

\textbf{Program status, revised negatives and reopened reports
(manuscript Section 5.3).} A merchant's network program status can also
change that meaning: under one program, a merchant loses the liability
shift on domestic authenticated payments until it exits the program
{[}\protect\hyperlink{ref-Stripe2026MonitoringPrograms}{6}{]}. The
records keep that status as it stood at decision time, because the guide
does not say whether the loss reaches payments authenticated before the
merchant entered the program. A later program entry is recorded as a
rule event, and the liability status of payments authenticated before it
remains unknown. A negative read from a matured absence cites the rule
version and read time under which it matured and is revised when a later
transition reopens the channel. A report's resolution is not always
final either: in rare cases an issuer files a further dispute on the
same payment after acquiring new information, or a lost dispute later
becomes a late win {[}\protect\hyperlink{ref-Stripe2026Disputes}{7}{]}.

\textbf{Bank-debit and other rails (manuscript Section 7.1).} A SEPA
Direct Debit dispute within eight weeks of the debit is honored with no
questions asked, and such disputes are final, with no process for appeal
{[}\protect\hyperlink{ref-Stripe2026SEPADebit}{8}{]}. After eight weeks
and up to 13 months, a customer can dispute only a debit considered
unauthorized, and the customer's bank decides; the window that applies
therefore depends on the claim, and an absence at eight weeks is not a
negative for unauthorized-debit fraud
{[}\protect\hyperlink{ref-Stripe2026SEPADebit}{8}{]}. Inside its window,
an ACH return is refunded without question, and the ACH Network has no
appeals process of its own
{[}\protect\hyperlink{ref-Stripe2026ACHTerms}{9}{]}. A SEPA dispute
honored with no questions asked is not by itself a fraud allegation. For
ACH returns, the provider's terms give 60 days for consumer accounts and
2 days for business accounts, counted from the customer's bank statement
{[}\protect\hyperlink{ref-Stripe2026ACHTerms}{9}{]}, while the Nacha
return deadlines that banks apply run from the settlement date of the
original debit {[}\protect\hyperlink{ref-LeadBank2026ACHReturns}{10}{]};
a channel record therefore names the rule from which it takes its clock
start. Banks may also accept late returns outside the window
{[}\protect\hyperlink{ref-Stripe2026ACHTerms}{9}{]}. The product scope
also lists instant push payments and stablecoin payments (S2); their
reporting channels, if any, follow each rail's own rules and windows,
which this paper did not review.

\needspace{4\baselineskip}

\textbf{Assessment store}

\begin{itemize}
\tightlist
\item
  Retain the model or reviewer identity and version. Keep unknown,
  not-assessed and failed-open states distinct from low risk.
  Failed-open means the assessment failed and activity continued.
\item
  Record when each assessment was made and frozen (its freeze time) and
  the as-of time of its inputs. For each model version, record its
  training and model-selection cutoff and the label rules it used. A
  label that became available after that cutoff cannot enter that
  version.
\item
  Never store an assessment as an observation of the conduct it
  assesses.
\item
  Link each assessment, whether a model, a detection or a reviewer
  produced it, and each detection or model version, to any agent output
  or released organizational-memory item that informed it. Record the
  writer's case or session and the agent outputs and released items open
  in it when the record was written.
\item
  For each detection or model version, link the validation records of
  responsibilities 16 and 17, with their version, date and result.
\end{itemize}

\needspace{4\baselineskip}

\textbf{Control-outcome store}

\begin{itemize}
\tightlist
\item
  \textbf{Decision:} retain the decision, who had authority to make it,
  the rule or model version, and the reason.
\item
  \textbf{Inputs and policy:} log the time of the decision and the
  inputs it used. Keep versions from which every other input available
  at that time can be reconstructed. Record the score, threshold or rule
  that produced the decision, and whether the policy was deterministic
  and, otherwise, the probability of each alternative. Identify any
  input that a decider, automated or human, used but that was not
  logged.
\item
  \textbf{Human decisions:} retain who decided, how the case was
  assigned and what the decider was shown.
\item
  \textbf{Dependencies:} link any agent output and any released
  organizational-memory item that informed the decision, so the decision
  can be revisited when that item expires or is corrected or retracted.
  Record the writer's case or session and the agent outputs and released
  items open in it when the record was written.
\item
  \textbf{Correctness:} record reviewed correctness when available, with
  how the case was selected for review (for example an appeal, a
  complaint, a sampling rule or review of every case) and, for a sampled
  review, its selection probability; otherwise record unknown.
\item
  \textbf{Error estimates:} keep separate links to each estimate, with
  its control version, cohort, outcome window, denominator, uncertainty
  and scope. The cohort is the fixed set of decisions the estimate
  covers.
\item
  \textbf{Channel effects:} record the decision's effect on each outcome
  channel as closed, delayed, changed in meaning, none or unknown. Fix
  these effects when the decision is recorded. Record later state
  transitions, including rule events, as separate records that never
  edit the decision.
\item
  \textbf{Restrictions:} distinguish allowing and blocking from
  restrictions, such as a hold that may later be released.
\item
  \textbf{Urgent containment:} for an action decided before quality
  review (manuscript Section 6.4), record that route, the impact class
  (none, degraded service or outage), any notification and restoration
  path, the review deadline and the later review's outcome, including
  any reversal.
\end{itemize}

\needspace{4\baselineskip}

\textbf{Structured threat-knowledge store}

\begin{itemize}
\tightlist
\item
  For inferred relationships, clusters, technique mappings and
  conclusions, retain confidence and the input and rule versions that
  produced them.
\item
  Record the writer as a correlation engine, agent proposal or person,
  and status as proposed, accepted or retracted. Record the person or
  accountable process that accepted or retracted the item, and when.
\item
  Retrieve an agent-proposed item only when it is marked as an
  unreleased agent inference. It never counts as corroboration.
\item
  Link an item that a person or correlation engine wrote to any agent
  output or released organizational-memory item that informed it. Record
  the writer's case or session and the agent outputs and released items
  open in it when the item was written. The item counts as dependent on
  those outputs and items.
\end{itemize}

\needspace{4\baselineskip}

\textbf{Organizational-memory store}

\begin{itemize}
\tightlist
\item
  Keep released items only, with the releaser, version hash, released
  claims and the confidence stated for each, scope, expiry, reviewer
  findings and originating draft.
\item
  Retain supersession and retraction links, and lineage from every cited
  observation to its root record. Record every later status record for a
  released item, such as a monitoring finding, an early expiry, a
  retraction or a withdrawal, with who or what wrote it and when.
\item
  For benign or exonerating items, keep a re-test date, a source of
  counterexamples independent of the writer, and the versions of the
  cited evidence and of any cluster the item concerns. A revision of
  that evidence or a change in that cluster's membership expires the
  item early (manuscript Section 6.5). Record each re-test, with its
  counterexamples' sources, its result and its date.
\item
  Read each item as context or a hypothesis, never as an observation.
\item
  Proposed lessons remain outside organizational memory until testing,
  quality review and human release. For each one, record its writer,
  proposal time, queue state, intended scope and version. Record each
  test against counterexamples, with the counterexamples' sources and
  the result. After release, record monitoring findings and any
  withdrawal, each with its writer and time, and the reason for a
  withdrawal (manuscript Section 7.3). The requirement for
  responsibility 23 in the functions file calls a withdrawal rollback;
  the two terms name the same record.
\end{itemize}

\needspace{4\baselineskip}

\textbf{External-knowledge store}

\begin{itemize}
\tightlist
\item
  Retain source, type, version, when the source observed and published
  the fact, validity window, license and retrieval time. The validity
  window is the period in which the fact holds. For a catalog of abuse
  techniques, the version is the fixed version that mappings cite.
\item
  Keep external knowledge distinct from platform observations.
\item
  Record the identifier attached to each released item shared outside
  the platform. When external content returns with that identifier, link
  it to that item. It counts as dependent on the item and on any agent
  output behind it, not as independent corroboration.
\end{itemize}

\needspace{4\baselineskip}

\textbf{Agent-and-approval ledger}

\begin{itemize}
\tightlist
\item
  For each run, record the human assignment or the standing task
  template and version that authorized it, with its permitted data,
  entities, period, current period and outputs. The current period
  bounds claims about current conduct.
\item
  For each tool call, preserve the provenance of inputs that reached the
  agent's language model.
\item
  For each approval, keep the approver, the target, task, tenant and
  tool it is bound to, its expiry, its reuse rule and any token issued.
  Keep hashes of the approved and executed payloads, and the
  tool-manifest version and hash in force at approval and at execution.
  The manifest is the declared list of the agent's tools.
\item
  Record memory writes and decisions on outbound traffic.
\item
  Record each agent output, such as an intelligence draft and each of
  its versions, a hunting or detection candidate or a recommendation,
  with its identifier, version hash, run and creation time. It is an
  observation of what the agent did, not evidence of the activity under
  investigation.
\item
  Retain model identity and version, instruction version and decoding
  settings.
\item
  Mark content that came from untrusted input as untrusted when it is
  read back. Distinguish copied untrusted content, agent output exposed
  to untrusted content and agent output without that exposure; exposure
  takes precedence.
\item
  Record each rule-checker result and quality-review finding with the
  draft's version hash, the rule version or reviewer, including the
  version of the rule that marks uncited returned records, and the
  draft's resulting status: returned, released or expired. For each
  quality review, record the materiality tier applied to each claim, the
  sampling rate, the cited records the reviewer examined and the random
  sample of unmarked uncited returned records it drew, whose size is
  fixed in advance for each claim type. Record when each draft version
  entered the review queue and when it would expire from it. For each
  review, record its start and end times.
\item
  Record each removal of an agent-output or released-item link from a
  label, control decision, detection, assessment or threat-knowledge
  item, with the person who removed it, that person's release authority,
  the time, the reason and the identifier of the removed output or item.
\item
  Store references to card data, identity documents and other regulated
  data, not copies (manuscript Section 5.5).
\end{itemize}

\needspace{4\baselineskip}

\textbf{Query coverage record}

\begin{itemize}
\tightlist
\item
  For each query run, retain its request: the task, run and question,
  the authorized requester, the tenant, purpose and data scope, the
  query version, the row, time and cost limits set, the expected
  observation and the required coverage.
\item
  For each query run, retain the coverage record the gateway attached:
  requested and evaluated entities with their tenant and identifier
  namespaces, requested and evaluated event intervals, as-of cutoff,
  source and schema versions, sampling, each source's event and
  ingestion watermarks, known gaps and retention horizon, access
  filtering, and any row, time or cost limit reached.
\item
  Keep the derived result status (complete, partial, truncated, denied
  or unavailable), the version of the rules that derived it and any
  scope-limited flag. The agent cannot write these fields.
\item
  For each query run, keep the identifiers of the records it returned,
  including those behind an approved aggregate, so that the checker can
  find returned records that a draft neither cites nor lists against a
  claim, attach those its fixed rule marks and summarize the rest
  (manuscript Section 6.4). Keep identifiers, not copies, under each
  record's retention and deletion rules (manuscript Section 5.5).
\item
  Link each tested hypothesis and each claim of absence to the query
  runs behind it. Record each hypothesis's status, including
  inconclusive and unassessable, apart from the query results. Record
  each incomplete test and each abstention, a decision not to conclude.
\end{itemize}

\needspace{4\baselineskip}

\textbf{Time semantics}

\begin{itemize}
\tightlist
\item
  Use timestamps precise enough to join a key's exposure to its use
  within minutes, with each timestamp's clock source.
\item
  For each payment, keep authorization, capture and processing times
  separately, and any service or event date that a channel rule uses as
  its clock start. Keep the authentication result and the reported
  liability status separately, each with its source, rule version and
  learned time. Keep any network program status that removes the
  liability shift as it stood at decision time. Record a later program
  entry as a rule event, and leave the liability status of payments
  authenticated before it unknown. An unknown liability status remains
  unknown.
\item
  Preserve label delay and what information was available as of each
  time.
\item
  Retain account age, dormancy and how long funds stay before moving on.
\item
  For each record class, keep each record until the longest applicable
  reporting window, counted from its channel's clock-start event, or the
  defined maturity period has ended, whichever ends later, plus the
  investigative lookback period. Follow the retention and deletion rules
  in manuscript Section 5.5. Record resulting gaps as coverage limits.
  Record each authorized hold with its scope and retention rule.
\end{itemize}

\needspace{4\baselineskip}

\textbf{Graph semantics}

\begin{itemize}
\tightlist
\item
  Link closed accounts to suspected successor accounts, and storefronts
  to merchants.
\item
  Link each outcome report to the transaction, payout or account
  application it concerns, with the basis for the link and whether it is
  verified. Attribute outcomes through that link to the platform
  account, connected account and end customer, so that a measure counts
  each affected transaction once. A report's links to the decisions that
  may have shaped it do not change this count. Keep each transaction's
  acquiring bank and the static part of its statement descriptor, so
  that a network's monitored unit can be reconstructed.
\item
  Link each wallet to every party it has allowed to spend its funds.
\end{itemize}

\FloatBarrier

\hypertarget{s5.-pinned-public-schemas}{%
\subsection{S5. Pinned public schemas}\label{s5.-pinned-public-schemas}}

\suppressfloats[t]

\href{SUPPLEMENT/S5_SCHEMA_VERSIONS.csv}{Schema versions} records the
filenames, hashes, versions and counts at repository commit
\texttt{\seqsplit{18fa2cc768024b47789dc6fa243acd48478acac2}}, dated 25
September 2026. All three files come from the public
\texttt{stripe/openapi} repository. The commit message reads ``Update
OpenAPI for 2026-09-30.endive'': the version strings name the API
release the files describe, which postdates the commit. The
\texttt{private\_preview} in file C's name does not indicate access to
private systems. Stripe publishes these files under the MIT License.
That license governs the schema-derived identifiers and locators in the
event-type and signal-locator files; the package's third-party notice
reproduces it.

The pinned public schemas supply event identifiers, creation time and
mode, but omit learned time and some actor, scope and lineage fields
{[}\protect\hyperlink{ref-StripeOpenAPI2026}{11}{]}. Stripe's
canonical-logging account describes one summary record per request in
each service, with identifiers and outcomes
{[}\protect\hyperlink{ref-Leach2019}{12}{]}. These are partial
precedents; the specification does not require vendor field names.

\Needspace{120pt}

File A is \texttt{latest/openapi.spec3.sdk.json}; B is
\texttt{preview/openapi.spec3.sdk.json}; C is
\texttt{preview/\allowbreak openapi.private\_preview.spec3.sdk.json}.
Paths are endpoint addresses. Resource schemas are component schemas
with an \texttt{x-resourceId}; they describe kinds of objects. Some
resource schemas in file C share an ID.\par\smallskip
{\footnotesize
\begin{tabular}{@{}
  >{\raggedright\arraybackslash}p{(\columnwidth - 8\tabcolsep) * \real{0.1100}}
  >{\raggedright\arraybackslash}p{(\columnwidth - 8\tabcolsep) * \real{0.2300}}
  >{\raggedleft\arraybackslash}p{(\columnwidth - 8\tabcolsep) * \real{0.1000}}
  >{\raggedleft\arraybackslash}p{(\columnwidth - 8\tabcolsep) * \real{0.3600}}
  >{\raggedleft\arraybackslash}p{(\columnwidth - 8\tabcolsep) * \real{0.2000}}@{}}
\toprule\noalign{}
\begin{minipage}[b]{\linewidth}\raggedright
SDK file
\end{minipage} & \begin{minipage}[b]{\linewidth}\raggedright
Version
\end{minipage} & \begin{minipage}[b]{\linewidth}\raggedleft
Paths
\end{minipage} & \begin{minipage}[b]{\linewidth}\raggedleft
Resource schemas / distinct IDs
\end{minipage} & \begin{minipage}[b]{\linewidth}\raggedleft
Event types
\end{minipage} \\
\midrule\noalign{}

A & 2026-09-30.endive & 421 & 175 / 175 & 508 \\
B & 2026-09-30.preview & 574 & 237 / 237 & 618 \\
C & 2026-09-30.preview & 800 & 329 / 327 & 781 \\

\bottomrule\noalign{}
\end{tabular}
}\par

\href{SUPPLEMENT/S5_EVENT_TYPES.csv}{Event types} lists the types
declared under \texttt{x-stripeEvent} in each file. Each count includes
216 thin-event declarations whose names add \texttt{v1.} to a snapshot
type in the same file. The counts do not also include the older
\texttt{event.type} list.
\href{SUPPLEMENT/S5_SIGNAL_LOCATORS.csv}{Signal locators} identifies
each counterpart's file and location using slash-separated field names
and list positions. Slashes within endpoint names and media types, such
as \texttt{application/json}, remain unchanged, so adjacent slashes can
occur, as in \texttt{/paths//v1/webhook\_endpoints}. The file's
\texttt{locator} column therefore does not hold standard JSON Pointers,
which would write each such slash as \texttt{\textasciitilde{}1}.

The signal-locator file lists only the 32 signals with partial
counterparts. For the seven with none identified, the classification
rests on a keyword search of all three files, so a counterpart under
other terms is not ruled out. That search covered eight signals:
\mbox{SIG-03}, \mbox{SIG-34}, \mbox{SIG-35}, \mbox{SIG-38},
\mbox{SIG-39}, \mbox{SIG-40}, \mbox{SIG-56} and \mbox{SIG-68}. Its terms
were \texttt{camera}, \texttt{emulator}, \texttt{inject},
\texttt{liveness}, \texttt{selfie}, \texttt{capture\_session},
\texttt{device}, \texttt{tamper}, \texttt{manipulat}, \texttt{spoof},
\texttt{presentation}, \texttt{jailbreak}, \texttt{rooted},
\texttt{virtual}, \texttt{agent}, \texttt{tool\_call}, \texttt{tool},
\texttt{mcp}, \texttt{llm}, \texttt{prompt}, \texttt{model\_context},
\texttt{manifest}, \texttt{tool\_version}, \texttt{hash},
\texttt{signature}, \texttt{signed\_request}, \texttt{web\_bot},
\texttt{bot}, \texttt{verified\_agent}, \texttt{attestation},
\texttt{memory}, \texttt{skill}, \texttt{egress}, \texttt{allowlist},
\texttt{proxy}, \texttt{outbound\_request}, \texttt{request\_log},
\texttt{user\_agent}, \texttt{ip\_address}, \texttt{tls}, \texttt{ja3},
\texttt{fingerprint}, \texttt{http\_status}, \texttt{api\_error},
\texttt{request\_id}, \texttt{api\_request}, \texttt{allowance},
\texttt{spender}, \texttt{permit}, \texttt{signing}, \texttt{erc20},
\texttt{smart\_contract}, \texttt{token\_approval}, \texttt{approve} and
\texttt{unlimited}. It searched event types, schema names, paths, keys,
enum values and descriptions. For \mbox{SIG-03}, identity-verification
outcomes such as \texttt{selfie\_manipulated} and the
\texttt{require\_live\_capture} setting record no device, camera,
emulator or injection context. For \mbox{SIG-34}, the
\texttt{agent\_details} fields name the agent behind a payment or token
and file C has MCP keys for claimable sandboxes, but neither records a
tool call or its input provenance. For \mbox{SIG-40}, an API key's
\texttt{ip\_allowlist} controls inbound use of the key, not egress. For
\mbox{SIG-56}, the event envelope's request ID and idempotency key, an
API error's \texttt{request\_log\_url} and aggregate API-error health
alerts give no per-request context. For \mbox{SIG-38}, the search found
the \texttt{agent\_details} fields, which the signal-locator file
records as its partial counterpart.

Every listed counterpart is partial. An event envelope (the wrapper
around an event) supplies only some required provenance. A request ID
does not supply request context, and a score does not establish the
behavior it scores. Likewise, current state does not supply a complete
history. Onramp-session objects (sessions for buying cryptocurrency) and
event-destination configurations (settings for event delivery)
illustrate this gap. The pinned schemas record part of their state
without complete matching event streams. Stripe's crypto documentation
describes a webhook for each onramp-session status change after
creation, but none on creation. Neither the event types that the three
files declare under \texttt{x-stripeEvent} nor their lists of
webhook-endpoint events name a crypto or onramp event.

\clearpage

\hypertarget{s6.-hypothetical-investigations-and-reproducible-arithmetic}{%
\subsection{S6. Hypothetical investigations and reproducible
arithmetic}\label{s6.-hypothetical-investigations-and-reproducible-arithmetic}}

\suppressfloats[t]

\begin{table}[t]
\textbf{Table S6.1. The two investigations at a glance.} Every
observation is hypothetical. Manuscript Sections 8.1 and 8.2 identify
the supporting research and how far its findings transfer to this
setting.\par\smallskip
{\small
\begin{tabular}{@{}
  >{\raggedright\arraybackslash}p{(\columnwidth - 4\tabcolsep) * \real{0.3000}}
  >{\raggedright\arraybackslash}p{(\columnwidth - 4\tabcolsep) * \real{0.3400}}
  >{\raggedright\arraybackslash}p{(\columnwidth - 4\tabcolsep) * \real{0.3600}}@{}}
\toprule\noalign{}
\begin{minipage}[b]{\linewidth}\raggedright
Element
\end{minipage} & \begin{minipage}[b]{\linewidth}\raggedright
Trace A: exposure and possible payout redirection
\end{minipage} & \begin{minipage}[b]{\linewidth}\raggedright
Trace B: storefront links and possible dispute evasion
\end{minipage} \\
\midrule\noalign{}

Trigger & An external observation of an exposed key & A platform risk
assessment and a rejected application \\
First questions & Did use change after exposure? What requests and
account changes followed? & Which storefronts share templates,
descriptors or merchant identifiers? \\
Main analysis & Join events by time; examine the history and reuse of
payout destinations & Examine hard and soft links, and outcomes linked
to each transaction, across accounts \\
Benign alternatives & Developer testing from a new network or an
internal scanner; a business with several legitimate accounts, or an
intermediary receiving funds for several clients & A legitimate
multi-brand business, shared template vendor or aggregator \\
Evidence limit & Payout redirection is reasoned about by analogy with
payroll diversion; timing premises come from Microsoft's report of
follow-on business email compromise payment fraud soon after session
theft and a deliberate key-exposure experiment & Complaint allegations
are not findings; resolving pre-dispute alerts can change chargeback
ratios \\
What limited evidence cannot show & A partial baseline cannot show that
there was no earlier misuse & Neither soft links nor a hard link from a
verified identifier by itself establishes common control, intent or who
is behind the activity \\
Proposed outputs & Hunting candidate; detection candidate with an
evaluation plan; recommendations to consider rotating or revoking the
credential and a payout pause or hold; intelligence draft & Cluster
intelligence assessment; recommendation to review the newly onboarded
merchant accounts linked to the storefronts; detection candidate built
on hard links between storefronts and accounts \\
Who acts & Authorized humans or accountable operational systems decide.
The urgent exception covers a reversible payout pause or hold; rotation
or revocation follows normal review & Authorized humans or accountable
operational systems decide whether to reject or terminate accounts; the
agent does not \\

\bottomrule\noalign{}
\end{tabular}
}
\end{table}

Table S6.1 summarizes the two investigations. The
\href{SUPPLEMENT/S6_SYNTHETIC_COHORT.csv}{synthetic cohort} contains 400
completed transactions, 100 in each of four non-overlapping account
cohorts. Each cohort fixes one account's transactions using the same
definition and outcome window. These synthetic rows make the counting
rules reproducible. They contain no report-arrival dates, so the ratios
assume that every report on the four cohorts has arrived. They cannot
show reporting delay.

Once each report is linked to its transaction, each of the four
accounts, A to D, has two distinct transactions with chargebacks, so its
chargeback ratio is \mbox{\texttt{2/100\ =\ 2\%}}. Pooling the same
measure across the four accounts gives \mbox{\texttt{8/400\ =\ 2\%}}.

Consider non-overlapping cohorts with the same outcome definition and
outcome window. Let \(T_i > 0\) be account \(i\)'s transaction count and
\(D_i\) its count of distinct transactions with chargebacks. Then

\[
r = \frac{\sum_i D_i}{\sum_i T_i} = \sum_i w_i r_i,
\qquad r_i = \frac{D_i}{T_i},
\qquad w_i = \frac{T_i}{\sum_j T_j}.
\]

The weights are nonnegative and sum to one, so the pooled ratio lies
between the smallest and largest account ratios.

Fourteen transactions have a resolved pre-dispute alert, including two
of the eight transactions with chargebacks. Counting each transaction
once gives \mbox{\texttt{8\ +\ 14\ −\ 2\ =\ 20}}. The separate
chargeback-or-alert-resolution measure is therefore
\mbox{\texttt{20/400\ =\ 5\%}}. Adding all flags would count two
transactions twice and wrongly give \mbox{\texttt{22/400\ =\ 5.5\%}}.
This broader measure is neither a fraud rate nor the chargeback ratio
above. The example assumes no dispute-monitoring threshold for either
measure. The broader measure is \texttt{8/100} for A and B, and
\texttt{2/100} for C and D; the pooled 5 percent lies between them.

The chargebacks on D001 and D002 are recorded at account-X, whose own
transactions lie outside all four cohorts. Both transactions remain in
D's cohort and count once; account-X adds no transactions to the
denominator. Counting reports where recorded shows D at \texttt{0/100}
and assigns both reports to account-X. Linking reports to transactions
restores D to \texttt{2/100}, the same as the other accounts. The
mismatch needs explanation but does not prove concealment. Legitimate
merchant-of-record or facilitator arrangements can also record disputes
on another account.

Both measures exclude reports on transactions outside the four cohorts.
A missing transaction link leaves a measure that depends on it
incomplete. Any change to report links or outcome definitions must be
stated explicitly. If it applies alike to every cohort, the
weighted-average formula still holds. The exact 95 percent intervals in
manuscript Section 8.2 are Clopper-Pearson intervals, and its p-value
comes from Fisher's two-sided exact test of 8/100 against 2/100.

Trace A uses a separate hypothetical timeline (manuscript Sections 8.1
and 8.3). The agent requests seven days of history before the key's
exposure; only the final two hours are available. A detection already
running on the platform freezes its first compromise assessment at +30
minutes after publication of the key. The destination-reuse link becomes
known at +1 day, although it was effective before exposure. The
+30-minute assessment cannot use that link, and the missing history
leaves the onset of any misuse inconclusive.

\url{CHECK_CANDIDATE.py} checks these calculations and excludes the
+1-day destination-reuse link from the +30-minute assessment. It also
checks identifiers, counts, feature lineage, citation-reference
agreement and links between files. \url{REPRODUCE.md} gives build and
verification instructions. These are checks of the supplied artifacts
and arithmetic, not an evaluation of the proposed system.

Manuscript Table 6 uses eleven signals: \mbox{SIG-04}, \mbox{SIG-15},
\mbox{SIG-16}, \mbox{SIG-18}, \mbox{SIG-20}, \mbox{SIG-24},
\mbox{SIG-25}, \mbox{SIG-32}, \mbox{SIG-56}, \mbox{SIG-58} and
\mbox{SIG-71}. Their complete lineage inputs and the storage
requirements are also needed. In particular, the denominator requires
the eligible completed-transaction cohort from \mbox{SIG-27}, including
transactions without reports.

\FloatBarrier

\hypertarget{s7.-evaluation-conditions}{%
\subsection{S7. Evaluation conditions}\label{s7.-evaluation-conditions}}

\suppressfloats[t]

These conditions explain the scope and comparisons behind the tests in
S7.1. E8's release-gate test appears there alongside the other eight; S8
adds practitioner judgments of availability and usefulness.

\textbf{E1. Do coverage and learned-time records reduce unsupported
conclusions?} The confirmatory test in S7.1 concerns cases with
incomplete history or links learned late, with a stratum of cases with
complete history for its non-inferiority measure. A broader timing study
would compare minute-level joins with coarser timing across a realistic
incident population. It would measure exposure-to-use and
session-issuance-to-use intervals, detection or triage gains, collection
costs and false alarms from benign changes. Microsoft's report of
payment fraud attempts soon after session theft and the deliberate
key-exposure experiment motivate precise timing; neither establishes its
typical benefit.

\textbf{E2. When can missing outcomes be estimated?} Define the outcome
channel, population and assumptions before selecting an estimator.
Simulation can test behavior under known assumptions, including cases
without overlap, where no comparable cases received another decision. It
cannot verify those assumptions or recover unknown real outcomes. Later
observational comparisons must separate populations with supported
overlap from those whose outcomes cannot be estimated. Test sensitivity
to late, mistaken and deliberately false reports. Allowing attempts that
a policy would otherwise block requires separate ethical, human and
operational authorization. The motivating theoretical preprint proposes
methods; this paper claims no recovery of missing outcomes on platform
data {[}\protect\hyperlink{ref-Dhama2026Causal}{13}{]}.

\textbf{Examining assumptions is not validating an estimator (manuscript
Section 7.1).} Three target quantities illustrate the distinction:

\begin{itemize}
\tightlist
\item
  \textbf{Candidate-policy value:} the average outcome a proposed
  decline policy would produce on a stated measure. Estimation from
  another policy's logs needs decision inputs, probabilities and overlap
  for the decisions the candidate would take.
\item
  \textbf{The chargeback rate if all cases were allowed:} this needs
  comparable cases with some chance of being allowed. A class always
  blocked supplies no allowed comparison.
\item
  \textbf{A control's error rate:} this needs a defined cohort, outcome
  window and denominator, with correctness reviewed for every case or a
  random sample with recorded selection probabilities. Appeals alone do
  not establish a cohort rate.
\end{itemize}

\textbf{E3. Do report linkage and cluster monitoring reveal evasion
earlier?} Compare the network's monthly rule, its ratio, its minimum
count and any minimum value, applied to each account as reports are
recorded, with the same rule applied to each cluster after linking
reports to transactions. In S7.1, each mention of the network's ratio
for E3 means this complete rule, including its minimum count and any
minimum value. A third arm applies the rule to each account after
linking, separating linkage from clustering. When every account has
captured payments in the denominator period, a pooled ratio cannot
exceed the highest account ratio, so a lead from clustering can come
only from a condition on totals, such as the minimum count or, where the
network sets one, a minimum value. An account with no captured payments
in the rule's denominator period has no defined ratio when it receives
reports, and the cluster ratio can then exceed every defined account
ratio. Report separately any cluster flag that depends on reports
received by an account with no captured payments in that period, such as
a closed account or, where the denominator is the previous month's
captured payments, an account in its first month. Keep the eligible
transactions, outcome definition and reporting periods consistent. Apply
the account arms to platform and connected accounts and, where a network
monitors a different unit, such as Visa's static statement descriptor
and acquiring bank, also to that unit. Evaluate broader measures,
including chargeback-or-alert-resolution, separately. Apply each
program's exclusions as the network's documentation states them, such as
VAMP's exclusion of fraud reports that qualify for Compelling Evidence
3.0. State whose report view each arm uses: the programs also count
disputes that the provider handles on a business's behalf and that do
not appear in that business's data
{[}\protect\hyperlink{ref-Stripe2026MonitoringPrograms}{6}{]}. Report
the share of reports linked to transactions, the share whose link is
unverified and how duplicate reports were removed. Measure false flags
on legitimate multi-brand businesses. Report erroneous links between
unrelated accounts sharing IP addresses, such as through a common proxy,
as a secondary measure; S7.1's failure criteria do not include it
{[}\protect\hyperlink{ref-Liu2018GEM}{14}{]}.

\textbf{E4. Can the records measure harm to legitimate users?} The
feasibility test in S7.1 asks whether control changes have computable
harm rates and intervals. It precedes E1 to E3 and E9. A broader
evaluation would compare detection gains with failed attempts and other
costs to legitimate users. The identity-verification and login-challenge
studies report user outcomes under their own procedures
{[}\protect\hyperlink{ref-Fatima2024}{15},
\protect\hyperlink{ref-Doerfler2019}{16}{]}. Their rates cannot
determine whether a proposed payment control is worthwhile. Rates from
the source studies require justification before transfer to this
platform.

\textbf{E5. Does the agent improve the investigation?} Compare the
governed agent with well-equipped human investigators and suitable
decision rules, including simple thresholds where appropriate. Sharma's
agent underperformed its underlying classifier's threshold on 60
synthetic cases
{[}\protect\hyperlink{ref-Sharma2026AuditableFraud}{17}{]}. That limited
result motivates testing the agent's added value. Give each approach the
same information, frozen at an as-of cutoff. Report case base rates and
judge outcomes without knowing which approach produced each finding. Set
error, abstention and escalation rules in advance. Measure correct
findings, errors, abstentions, human review burden, cost and latency,
and include adaptive adversaries. The study records reviewer time for
investigation, tool or agent operation, checking, release and rework on
the same basis in both human arms. The storage contract holds these
times only for work recorded in the agent-and-approval ledger.

\textbf{E6. Do the agent's safeguards hold under adaptive attack?} Test
prompt injection, approval rebinding (reusing approval for an altered
operation), memory poisoning, limits on outbound data, planted omissions
of retrieved records that count against a claim, and steering the agent
away from queries that would return such records. Let adversaries adapt
to the defense. Measure security failures and the effect of safeguards
on investigations completed. Prior adaptive attacks motivate this test
but do not establish that these safeguards work
{[}\protect\hyperlink{ref-Nasr2025}{18}{]}. One proposed defense
separates control flow, the steps an agent chooses, from untrusted data
{[}\protect\hyperlink{ref-Debenedetti2025CaMeL}{19}{]}. It is a
candidate for evaluation under its stated assumptions. Adopting it would
change the single hypothesis loop in manuscript Figure 3, which
currently makes no such separation.

\textbf{E7. Can lessons be reused without spreading mistakes?} Measure
time to useful reuse, negative transfer (a lesson degrading later
investigations) and poisoning. Vary validation strictness, release scope
and withdrawal conditions. General-purpose agent research supplies
methods and evidence of failures, not payment-platform results
{[}\protect\hyperlink{ref-Dash2026}{20},
\protect\hyperlink{ref-Xiong2026MemoryManagement}{21}{]}.

\textbf{E9. Do channel states reduce error against a mature holdout
reference?} Use decisions the platform already logs, with their
disputes, early fraud warnings, returns and other reports. For a fixed
cohort, build labels once under the platform's current label rule,
including its maturity period and its treatment of blocked or reversed
items, and once under the channel-state and maturity rules (manuscript
Section 5.3). If the platform has no such rule, use for each rail one
maturity period equal to that rail's standard window, with blocked
attempts excluded. Compare the resulting labels, an existing control's
measured error rates and the ranking of two candidate rules. On
decisions with a randomized holdout, compare the absolute error of the
same evaluation quantity under both readings against the holdout
reference after the longest defined reporting window has closed, with
report truth reviewed on a random sample. The target is the selected
metric on matured reported outcomes after sampled truth review. Only the
reduction in reference error, with joint uncertainty as specified in
S7.1, counts toward E9's pass rule. This comparison does not by itself
estimate statistical bias or reveal fraud that was never reported.
Report, by channel, how often the records determine an affected item's
channel state, clock start and window rather than recording them as
unknown. Record and report any rule event that reopens a window during
follow-up. The motivating theoretical preprint models reporting and
maturity for one channel
{[}\protect\hyperlink{ref-Dhama2026Causal}{13}{]}; E9 asks whether
recording several channels' states matters on platform data.

Record the provenance of synthetic, replayed and emulated data when
created. Agent-driven adversary emulation would also require the
authorization and safeguards in manuscript Section 6.5; this paper
performs and authorizes none.

\hypertarget{s7.1-pre-specified-hypotheses-and-criteria}{%
\subsubsection{S7.1 Pre-specified hypotheses and
criteria}\label{s7.1-pre-specified-hypotheses-and-criteria}}

For each question in manuscript Section 9.1, this section proposes a
hypothesis, data design, primary measure and result that would count
against the specification. E2's interval settings remain
deployment-specific choices, to be fixed before any data are examined.
Where a question in Section 9.1 is broader, the test below addresses the
part that bears on the three record requirements or, for E5, on the
agent's added value. Thresholds and sample sizes are proposed defaults.
E1's non-inferiority margin, E3's false-flag margin and E8's tolerable
and assumed release rates are named parameters that a platform fixes and
records before data collection, and E9's negligible reduction in
reference error is one that it fixes before any data are examined; this
section does not set them. A platform may change them before it starts,
if it records the change before examining any data. No data have been
collected, and whoever runs the plan can register it with a time stamp.

Four rules apply to all nine. First, each evaluation freezes the
information each method may use at an as-of cutoff and separates
development from test data. Second, it states how true outcomes are
decided, by whom, and blind to which method produced each finding. Its
error and abstention rules are set in advance, and it states base rates.
Third, where ground truth comes from labels shaped by control decisions,
it says how that bias is avoided. Otherwise it reports relative
performance among detected cases only. Fourth, records from tests,
replays and emulations are marked when created and kept out of
production labels and lessons.

Each evaluation is judged separately, without correction for multiple
tests, because support requires every confirmatory evaluation to pass.
Each confirmatory evaluation targets 93 percent power: the chance of
passing its complete rule under stated effect assumptions. If each
reaches that target, all four pass with probability at least 72 percent
by the union bound, \(1-4\times0.07\), or about 75 percent if their
results are independent. Sizing must account for every pass criterion
and case or reviewer clustering. The minimum counts below are not power
calculations; the sizing assumptions must be recorded before main-study
data are collected. E1, E3, E8 and E9 are confirmatory and the rest are
exploratory. Support in a setting requires every confirmatory evaluation
to pass its complete rule. Each criterion has three possible results.
Support: its pass rule is met. Evidence against: a bound at the same
confidence level, taken from the other side, lies on the failing side of
its threshold; for example, E3 shows excess false flags if the lower
one-sided 95 percent bound on the increase is at or above the margin.
Inconclusive: neither holds, or the target cannot be identified from the
data. A failure criterion below means only that support was not
established; the interval and the identification limits decide which of
the other two results applies. No result by itself shows that the
underlying idea is wrong. Investigators cannot be blind to the added
records, so the raters, not the investigators, are blind. A lower or
upper 95 percent bound is an end of a two-sided 95 percent interval
unless it is marked one-sided.

\hypertarget{e1-tests-the-second-requirement-confirmatory.}{%
\paragraph{E1 (tests the second requirement;
confirmatory).}\label{e1-tests-the-second-requirement-confirmatory.}}

\textbf{Hypothesis:} coverage and learned-time records reduce the share
of unsupported conclusions on cases with partial history or with links
learned late, without reducing the correct, supported conclusions
reached per case.

\textbf{Data and design:} closed real cases of that kind, each worked by
different investigators with and without those records, and rated by two
blind raters. Investigators are assigned to the two conditions at
random, balanced so that each works cases in both conditions and none
works the same case twice. Raters score conclusions as single
propositions under a rubric fixed in advance: they split compound
conclusions into their propositions and merge restatements of one
proposition, so that each distinct proposition is scored once however it
is phrased or grouped. A stratum of closed cases with complete history
is worked and rated the same way. Who investigates is fixed before the
study starts: people, the agent, or both as separate arms, each analyzed
on its own. If both arms are run, E1 fails if either arm meets a failure
criterion, and the power target applies to both arms together.
Conclusions reach the raters in a common format without coverage or
learned-time field names. Each rater records which arm they believe
produced each conclusion, so that blinding can be checked. A 30-case
pilot, excluded from the main study, estimates the level without the
records for sizing. The main study will be sized for 93 percent power
for its complete pass rule, assuming the share truly falls by half and
correct, supported conclusions per case do not fall. Intervals allow for
several conclusions per case.

\textbf{Primary measures:} the share of conclusions rated unsupported by
the evidence then available and, co-primary, the number of correct,
supported conclusions per case, rated by the same blind raters against
the closed case's resolution. Abstentions per case are reported for each
arm and stratum. Because the records can lower the share of unsupported
conclusions simply by leading investigators to abstain, E1 also reports
supported conclusions and abstentions per case.

\textbf{Failure criteria:} Fails if, in the cases with partial history
or links learned late, the share falls by less than a third of the level
without the records, or the 95 percent interval on the paired difference
includes zero. It also fails if, in those cases or in the
complete-history stratum, the lower 95 percent bound on the difference
in correct, supported conclusions per case (with the records minus
without) is at or below minus the non-inferiority margin.

\hypertarget{e2-tests-the-first-requirement-exploratory.}{%
\paragraph{E2 (tests the first requirement;
exploratory).}\label{e2-tests-the-first-requirement-exploratory.}}

\textbf{Hypothesis:} the control-decision record is complete enough to
estimate outcomes under other decisions (manuscript Section 7.1).

\textbf{Data and design:} a simulation with known effects first, then
the last quarter of decisions where the platform already randomizes
traffic under its own authorization. Record that quarter's start and end
dates before examining the data. The target is matured reported
outcomes: holdout rates count the reports received once each channel's
window has closed, and report truth is reviewed on a random sample. A
usable probability is one logged with the feature values and model
version that produced it. The deployment chooses the intended coverage
defined under the primary measures; it must be at least 95 percent.

\textbf{Primary measures:} the share of randomized decisions logged with
a usable probability, feature values and model version. Also, the share
of overlap populations whose holdout rates fall outside the predicted
interval. Each interval predicts the observed outcome rate in a held-out
population, including its sampling variation. Before examining any data,
specify the estimator, interval construction, intended coverage and
overlap populations. Intended coverage is the planned share of repeated
holdout rates the intervals contain. These deployment-specific settings
are not selected here.

\textbf{Failure criteria:} Fails if fewer than 95 percent of randomized
decisions carry a usable probability, or more overlap populations fall
outside their predicted intervals than the intended coverage plausibly
allows: a one-sided binomial test, treating the populations as
independent, gives p below 0.05.

\hypertarget{e3-tests-report-linkage-which-the-first-requirement-depends-on-together-with-cluster-monitoring-confirmatory.}{%
\paragraph{E3 (tests report linkage, which the first requirement depends
on, together with cluster monitoring;
confirmatory).}\label{e3-tests-report-linkage-which-the-first-requirement-depends-on-together-with-cluster-monitoring-confirmatory.}}

\textbf{Hypothesis:} linking each report to its transaction, and
computing the network's monthly ratio for each linked cluster, flags
evasion earlier than the network's ratio for each account.

\textbf{Data and design:} at least 30 linked-account clusters with later
confirmed outcomes, channels kept apart, and at least 100 legitimate
multi-brand businesses from onboarding. The pass rule is judged only on
clusters selected by a confirmation route documented as independent of
either monitoring rule; clusters from other routes are reported
separately. Adjudicators blind to both rules' outputs fix the onset of
evasion for each cluster from case evidence. Each program is applied as
the network's own documentation states it, with the provider's guide
used only as a secondary summary. Time to flag is analyzed as paired,
right-censored time-to-event data: a cluster that a method has not
flagged by the end of follow-up is censored there for that method. A
third arm computes the network's ratio for each account after linking.
It separates the effect of linking from that of pooling accounts into a
cluster. The cluster and third arms count only reports with verified
links, such as a match on the network's transaction identifier. Flags
that need reports linked otherwise, such as by card number and amount
alone, are reported separately. The failing result applies to the stated
contrast, and the third arm attributes any lead. Estimate the difference
in false-flag probabilities using the same legitimate businesses under
both rules, preserving business-level pairing and any additional
clustering. Prespecify the interval method and size the study for the
complete lead-time and false-flag pass rule at the 93 percent power
target, given the expected share of verified report links.

\textbf{Primary measures:} days from the adjudicated onset of evasion to
the flag, counting a flag on or before the onset as day 0, and false
flags. The lead time is the difference in restricted mean time to flag,
account arm minus cluster arm, up to a horizon fixed before any data are
examined.

\textbf{Failure criteria:} Fails if the lower bound of the 95 percent
interval on the lead time, estimated from the paired censored analysis,
is zero or less. It also fails unless the one-sided 95 percent upper
confidence bound on the increase in the probability that a legitimate
multi-brand business is flagged, cluster arm minus account arm, is below
the false-flag margin, proposed as 5 percentage points. The interval
preserves business-level pairing and any additional clustering. An
interval that crosses the margin is inconclusive about this guardrail;
an observed difference below the margin alone does not establish
support.

\hypertarget{e4-tests-the-first-requirement-exploratory.}{%
\paragraph{E4 (tests the first requirement;
exploratory).}\label{e4-tests-the-first-requirement-exploratory.}}

\textbf{Hypothesis:} for each control change (a new rule, a threshold
move or a model version), the records support a rate of harm to
legitimate users with a defined denominator.

\textbf{Data and design:} a sample of control changes, run before E1 to
E3 and E9 as a feasibility check.

\textbf{Primary measure:} the share of sampled changes with a computable
rate and interval.

\textbf{Failure criteria:} Fails if fewer than 80 percent of sampled
changes have one.

\hypertarget{e5-tests-the-agents-added-value-exploratory.}{%
\paragraph{E5 (tests the agent's added value;
exploratory).}\label{e5-tests-the-agents-added-value-exploratory.}}

\textbf{Hypothesis:} the agent with human review matches or improves
confirmed findings per reviewer hour compared with a person using tools.
Case error rates are no more than 2 percentage points higher than either
comparator arm's.

\textbf{Data and design:} at least 700 frozen cases and three arms:
threshold rules, a person with tools, and a person with the agent.
Errors are counted per case. Both human arms see the same records. In
E5, reviewer time covers human investigation, tool or agent operation,
checking, release and rework. Measure these activities on the same basis
in both human arms. Include work on cases that produce no correct
finding. As a supplementary analysis, set the threshold-rules arm to act
on as many cases as the person-with-agent arm, reporting its results by
depth of prior history and for fraud bursts defined in advance.

\textbf{Primary measures:} correct findings, errors, abstentions,
reviewer minutes per correct finding, cost and latency.

\textbf{Failure criteria:} Fails if errors in the agent arm exceed those
of the human arm or the threshold-rules arm by more than 2 percentage
points, or reviewer minutes per correct finding are higher. For minutes,
higher means that the 95 percent interval on the agent arm's ratio to
the human arm lies wholly above one. Avoiding a failure criterion does
not by itself establish equal error rates or improved efficiency.

\hypertarget{e6-tests-the-third-requirement-exploratory.}{%
\paragraph{E6 (tests the third requirement;
exploratory).}\label{e6-tests-the-third-requirement-exploratory.}}

\textbf{Hypothesis:} the safeguards prevent attacker-influenced content
from gaining unauthorized control over tool calls, outbound traffic,
released claims and the evidentiary status of records.

\textbf{Data and design:} a red team that adapts to the defenses from
the four attacker positions in manuscript Section 6.5, each with its
attempts, budget and model access stated. Content is planted in every
input the agent reads, and every field the agent can write is checked
for it. Planted content includes attempts to make the agent omit
retrieved records that count against a claim and attempts to steer it
away from queries that would return such records.

\textbf{Primary measures, per position:} attempts that cause an
unauthorized tool call, transfer or memory write; released material
claims that a blind adjudicator judges false; and missing links back to
agent outputs.

\textbf{Failure criteria:} Fails if any write to released organizational
memory occurs without human release or any required link back is
missing. It also fails if a false claim about a named entity is
released, either clearing it or stating that activity is absent. Success
for each position is reported with its budget, without a pass mark.

For these criteria, a write is one that adds or changes a released
claim; monitoring, expiry and withdrawal records are status records.
Against the insider position, a release or link removal by that insider
is an authorized action. E6 reports whether the release record, the
logged link removal, expiry and correction detect and limit it, and the
failure criteria apply to releases and removals that position did not
make.

\hypertarget{e7-tests-the-third-requirement-exploratory.}{%
\paragraph{E7 (tests the third requirement;
exploratory).}\label{e7-tests-the-third-requirement-exploratory.}}

\textbf{Hypothesis:} re-testing, scoped release and withdrawal keep a
reused lesson from degrading later investigations without delaying
useful lessons.

\textbf{Data and design:} lessons released under varied strictness,
planted poisoned lessons, and an expiry schedule disclosed to the red
team.

\textbf{Primary measures:} the rate of degraded later investigations,
the time to reuse, and the planted lessons released and still active.

\textbf{Failure criteria:} Fails if a planted poisoned lesson is
released and still active after its expiry or re-test date. It also
fails if the median time to reuse exceeds 14 days.

\hypertarget{e8-tests-the-third-requirement-confirmatory.}{%
\paragraph{E8 (tests the third requirement;
confirmatory).}\label{e8-tests-the-third-requirement-confirmatory.}}

\textbf{Hypothesis:} reviewers and releasers catch planted wrong drafts
when cited records are checked at the rate and materiality tiers planned
for deployment and they do not know which drafts are planted. Claims
about named entities are checked in full, and a person releases every
draft.

\textbf{Data and design:} planted drafts, interleaved with production
drafts at realistic volume and at a stated prevalence close to the
expected production error rate, spread across reviewers, with intervals
that allow for clustering by reviewer. The realistic-prevalence arm must
reach the required numbers of planted drafts and planted false
exonerations by itself, because both failure criteria apply to it. A
higher-prevalence arm may be added only for secondary estimates by error
type; it does not count toward the failure criteria, and both arms are
reported. Reviewers and releasers are not told which drafts are planted.
Planted drafts concern synthetic entities or closed cases whose cited
records are test-marked and whose named entities have no account action
pending. They carry the test environment mark when created;
reviewer-facing views hide it, while the stores enforce their exclusion
from production labels, intelligence and lessons. The gateway refuses
any control request or urgent-containment routing that originates from a
test-marked draft, and it tells the decider that the draft is planted
before any action. Released planted drafts are withdrawn after scoring.
Planted errors include false exonerations, drafts that wrongly clear a
named entity; omissions of retrieved records that count against a claim,
including records that the checker's fixed rule marks and records it
leaves to the sample (manuscript Section 6.4); and drafts whose only
support for a material claim is a released item or a record linked to
one. The mix of planted error types is fixed before the study starts.
Some planted errors come from the agent's own earlier runs and the rest
are written independently. A lower catch rate on the agent's own errors
would indicate shared errors, although the two sets may also differ in
difficulty. As specified, E6 does not separate errors the reviewer
shares with the agent from other misses. The number of planted false
exonerations is sized for 93 percent power to show a release rate below
the tolerable rate, given an assumed true release rate stated in
advance, and the number of planted drafts is sized so that the complete
rule, with the catch-rate criterion, reaches that power. The following
counts illustrate the false-exoneration endpoint under independent
Bernoulli trials with a common release probability and a fixed sample
size. They do not account for reviewer or case dependence or the
separate catch-rate criterion; the actual study must be sized for its
complete rule under the prespecified clustered design. With none
released, showing a rate below 1 percent at one-sided 95 percent
confidence needs at least 299, the smallest \(n\) with \(0.99^n<0.05\).
If the true rate is 0.1 percent, 93 percent power to show a rate below 1
percent needs 628, the smallest such number. With 628, two or fewer
releases show it, because at a true 1 percent rate that many occur with
probability 4.98 percent; at 0.1 percent they occur with probability
97.4 percent. If planted false exonerations made up 5 percent of the
drafts reaching review, 628 of them would take 12,560 drafts, about 628
days at the 20 drafts a day assumed in manuscript Section 6.4. The plan
states the tolerable rate, the assumed rate, the prevalence and the
implied duration together.

\textbf{Primary measures:} the catch rate, overall and by error type
(false exonerations, other wrong claims, omissions of retrieved contrary
records, and support resting only on released items or records linked to
them); the release rate of planted false exonerations; reviewer minutes
per released draft; and how often the recorded random sample of unmarked
uncited returned records contained a record that counted against the
claim. For the catch rate by reviewers and the releaser, a planted error
counts as caught only if one of them identifies it, by naming the wrong
claim or the missing record; a draft returned or left to expire without
that finding does not count. The checker's catch rate counts planted
errors that its fixed rules reject. Catch rates are reported for the
checker, reviewers and releaser, for each prevalence arm, and for errors
from the agent's own runs and independently written errors.

\textbf{Failure criteria:} Fails if the lower 95 percent bound on the
catch rate by reviewers and the releaser is below 80 percent, counting
planted errors that the rule checker did not reject, or the one-sided 95
percent upper bound on the release rate of planted false exonerations
exceeds the tolerable rate stated in advance. Both criteria apply to the
realistic-prevalence arm alone.

The release gate has a cost that E8 measures. If a system produced 20
drafts a day and a person needed 30 minutes to read each draft's cited
records, review would take 10 hours a day. The gate therefore moves
analyst work and may not remove it. At realistic volume the reviewer may
need to read a sample of each draft's cited records, tiered by
materiality. Claims about named entities are always checked in full, and
a person still releases every draft.

\hypertarget{e9-tests-the-first-requirements-channel-states-confirmatory.}{%
\paragraph{E9 (tests the first requirement's channel states;
confirmatory).}\label{e9-tests-the-first-requirements-channel-states-confirmatory.}}

\textbf{Hypothesis:} compared with the platform's current label rule,
channel-state and maturity records reduce the absolute error of an
evaluation quantity against a mature holdout reference for the same
prespecified target.

\textbf{Data and design:} two parts with one comparator. The comparator
is the platform's current label rule, with its existing maturity period
and its treatment of blocked or reversed items, recorded before any data
are examined. If the platform has no such rule, the comparator uses for
each rail one maturity period equal to that rail's standard window, with
blocked attempts excluded. (a) Historical reclassification: a fixed
cohort of payment attempts, payouts and account decisions, with the
decision logs, disputes, early fraud warnings, returns and other reports
the platform already holds. Labels are built twice, once under the
comparator and once under C1, which reads a missing report as a
provisional negative only where its channel matured open under its
original meaning. The cohort, channels and read times are recorded
before any data are examined. (b) Error against a randomized holdout
reference: decisions where the platform already randomizes traffic under
its own authorization, as in E2. At operational read times, an
evaluation quantity that the platform computes today, such as a
detection's precision or a control's false-positive rate, is computed
under both readings. Before any data are examined, the plan states how
items that either rule leaves unlabeled enter the quantity, whether
excluded, bounded or estimated under a stated delay model, and how
channels without a defined reporting window enter the holdout
comparison. The quantity, its population, denominator and read times,
and the method applied to unlabeled items are the same under both
readings, so that the comparison tests the records rather than a change
of estimator. Each reading is compared with a reference for the same
quantity and target population constructed from the holdout after the
longest defined reporting window has closed, with report truth reviewed
on a random sample. The target is the selected metric on matured
reported outcomes after that sampled truth review. Prespecify how the
truth-review sample contributes to the reference and its uncertainty.
Checking reports does not reveal fraud that was never reported; a target
of underlying fraud would require a separate ascertainment and
identification design. Prespecify a joint uncertainty procedure for the
two readings and the reference that preserves pairing, randomization,
truth-review sampling and relevant clustering. Rule events that reopen a
window during follow-up are recorded and reported. The main study will
be sized for 93 percent power for part (b), which carries the pass rule,
under stated effect assumptions.

\textbf{Primary measures:} (b) the reduction in absolute error against
the holdout reference, with its 95 percent interval. For comparator
estimate \(q_0\), C1 estimate \(q_1\) and holdout reference \(r\), the
reduction is \(D=|q_0-r|-|q_1-r|\); positive values favor C1. The
interval includes uncertainty in both readings and in the reference
under the prespecified joint procedure. Because \(D\) uses absolute
values, it is not smooth where \(q_1=r\), which is where C1 succeeds.
Standard delta-method and bootstrap intervals can miss their stated
coverage there, and the observed \(|q_1-r|\) tends to overstate C1's
error. Name the interval method before any data are examined, for
example a paired bootstrap that resamples randomization units within
truth-review strata, and size part (b) by simulating that same method.
This endpoint measures agreement with that reference. Statistical bias
instead concerns an expected estimate minus an identified target over a
specified sampling process and is not estimated merely by this realized
comparison. Reported but not counted as support: (a) the share of the
comparator's negatives that C1 classes as closed, delayed, changed in
meaning, immature or without a defined window. Items whose channel state
is unknown are excluded from that share and reported separately as a
coverage shortfall, with, by channel, the share of affected items whose
channel state, clock start and window can be determined rather than
recorded as unknown. Also reported: the reports that later arrive on the
comparator's negatives, and the change in an existing control's measured
error rates and in the ranking of two candidate rules.

\textbf{Failure criteria:} Fails if the lower 95 percent bound on the
reduction in reference error is at or below the negligible reduction, a
named parameter fixed before any data are examined. Measure (a) does not
count toward passing, because the share that C1 reclassifies follows
partly from its definitions and from the cohort's composition.

\FloatBarrier

\hypertarget{s8.-practitioner-review-protocol-not-yet-run}{%
\subsection{S8. Practitioner review protocol (not yet
run)}\label{s8.-practitioner-review-protocol-not-yet-run}}

\suppressfloats[t]

This proposed protocol assesses whether the specification fits
investigative practice. It measures availability and perceived
usefulness, not detection performance. Recruit six to ten fraud, abuse,
risk, trust-and-safety, data-engineering or privacy practitioners from
at least three organizations, with consent and required institutional
ethics review.

Give each participant the signal inventory (S3), responsibilities (S4)
and traces (S6), with and without the coverage, learned-time,
decision-link and release records. For each signal, ask whether it is
collected today (yes, partly, no), its typical retention and its
freshness (at decision time, in minutes, daily). Also ask about
usefulness in recent cases and the cost to an adversary of evading it.
For each responsibility, ask whether it exists and which role owns it.
For each trace, record the conclusions the participant would draw in
each condition.

Participants may also review a small set of planted wrong drafts as a
pilot for E8; such a pilot cannot reach E8's realistic volume and does
not test its criteria. Publish the protocol with ratings aggregated
across organizations and with free-text quotes stripped of identifying
detail, identifying signals rated missing, redundant or infeasible. No
published rating or quote may show which organization lacks a signal,
because that would disclose a control gap. No participant data have been
collected.

\FloatBarrier

\hypertarget{s9.-extended-comparison-with-prior-work}{%
\subsection{S9. Extended comparison with prior
work}\label{s9.-extended-comparison-with-prior-work}}

\suppressfloats[t]

Manuscript Section 2 develops the closest precedents for C1 to C3. This
section places them within research on payment decisions, investigation
and agent memory. The comparisons concern the cited designs and
findings, not every capability their systems might support. S1 records
source versions, publication status and supporting passages. Table S9.1
summarizes the nearest precedents for each requirement and what the
specification adds to them; the subsections that follow give the wider
context.

\begin{table}[!b]
\textbf{Table S9.1. The three requirements and their nearest
precedents.} The middle column states what the inspected sources
establish; the right column states what the specification adds for
cross-product investigation on a payment platform. The search was not
systematic (manuscript Section 9.2).\par\smallskip
{\footnotesize
\begin{tabular}{@{}
  >{\raggedright\arraybackslash}p{(\columnwidth - 4\tabcolsep) * \real{0.1500}}
  >{\raggedright\arraybackslash}p{(\columnwidth - 4\tabcolsep) * \real{0.5300}}
  >{\raggedright\arraybackslash}p{(\columnwidth - 4\tabcolsep) * \real{0.3200}}@{}}
\toprule\noalign{}
\begin{minipage}[b]{\linewidth}\raggedright
Requirement
\end{minipage} & \begin{minipage}[b]{\linewidth}\raggedright
Nearest precedents and what they establish
\end{minipage} & \begin{minipage}[b]{\linewidth}\raggedright
What the specification adds
\end{minipage} \\
\midrule\noalign{}

C1: outcomes linked to decisions & Selective-labels evaluation assumes
data in which each case carries the decision taken
{[}\protect\hyperlink{ref-Lakkaraju2017Selective}{22}{]}; learning from
logged decisions treats the decision and its probability as required
data {[}\protect\hyperlink{ref-Bottou2013Counterfactual}{23}{]};
decision provenance records the data flows around a decision
{[}\protect\hyperlink{ref-Singh2019DecisionProvenance}{24}{]}. In
consumer credit, reject inference faces the same problem for rejected
applicants; a review finds the common methods' assumptions mostly
implicit and no method uniformly better
{[}\protect\hyperlink{ref-Ehrhardt2021RejectInference}{25}{]}. Dhama's
missing-data model carries, for each transaction, the authorization
decision, whether fraud was reported and whether that report had arrived
by training time, for one reporting channel, and models label corruption
separately {[}\protect\hyperlink{ref-Dhama2026Causal}{13}{]}; a second
theoretical preprint by the same author bounds learning under delayed,
censored, corrupted and counterfactually missing feedback, treating each
as a distinct impairment
{[}\protect\hyperlink{ref-Dhama2026Limits}{26}{]}. Adyen evaluates
payment-time interventions offline from transactions logged under the
existing policy {[}\protect\hyperlink{ref-Egg2025AdyenOPE}{27}{]}; a
Meta preprint proposes estimators for the false positive rate of
blocking, which cannot be monitored directly once a user is blocked, and
tests them on simulated blocking policies
{[}\protect\hyperlink{ref-Tocker2022VirtualControl}{28}{]}; a
fraud-operations preprint treats automating an action as an
authorization decision because its labels are selective and delayed
{[}\protect\hyperlink{ref-Deng2026FCAC}{29}{]}; a decision-event schema
records facts learned after a decision, such as when ground truth
arrived, as linked records that reference the decision
{[}\protect\hyperlink{ref-Solozobov2026DecisionTrace}{30}{]}; the
coverage model in the C2 row also links records to decisions and
distinguishes refused execution branches from late reports
{[}\protect\hyperlink{ref-Wu2026EvidenceCoverage}{31}{]}; a 2015 talk
abstract describes adding randomness to production fraud scoring so that
a blocking model can be evaluated, and new models trained, although
blocked transactions' outcomes are not observed
{[}\protect\hyperlink{ref-Manapat2015Counterfactual}{32}{]}. & Several
outcome channels per decision, each with its own clock-start event and
reporting window; a record of what each decision closed, delayed or
changed in each channel, and a status for each channel at each read
time; outcome reports as separate appended records with typed links to
every decision that may have shaped them; separate fields for delay,
unobservability and report truth, so the evaluation states joint
assumptions; the recorded policy probability, inputs and case assignment
that those assumptions need (manuscript Section 7.1). A platform that
already joins decision logs to dispute records can use C1 to specify
what the join must preserve. \\

\bottomrule\noalign{}
\end{tabular}
}
\end{table}

\begin{table}[!tp]
\textbf{Table S9.1 (continued).}\par\smallskip
{\footnotesize
\begin{tabular}{@{}
  >{\raggedright\arraybackslash}p{(\columnwidth - 4\tabcolsep) * \real{0.1500}}
  >{\raggedright\arraybackslash}p{(\columnwidth - 4\tabcolsep) * \real{0.5300}}
  >{\raggedright\arraybackslash}p{(\columnwidth - 4\tabcolsep) * \real{0.3200}}@{}}
\toprule\noalign{}
\begin{minipage}[b]{\linewidth}\raggedright
Requirement
\end{minipage} & \begin{minipage}[b]{\linewidth}\raggedright
Nearest precedents and what they establish
\end{minipage} & \begin{minipage}[b]{\linewidth}\raggedright
What the specification adds
\end{minipage} \\
\midrule\noalign{}

C2: query coverage and learned time & Temporal databases separate valid
time from transaction time
{[}\protect\hyperlink{ref-Snodgrass1986Temporal}{33}{]}; completeness of
a query answer can be derived from statements about which parts of the
data are complete
{[}\protect\hyperlink{ref-Razniewski2011Completeness}{34}{]}; the
Dataflow model assumes a system processing unbounded data never knows
whether all data have arrived, treats watermarks as heuristic estimates
of progress, and revises results when late data arrive
{[}\protect\hyperlink{ref-Akidau2015Dataflow}{35}{]}. Bitemporal agent
memories keep both times
{[}\protect\hyperlink{ref-Wang2026TOKI}{36}--\protect\hyperlink{ref-Dan2026Anamnesis}{38}{]},
and Quipu gates writes to a knowledge graph while retaining validity and
recording times, with some policies requiring human approval bound to
the submitted evidence {[}\protect\hyperlink{ref-Brown2026Quipu}{39}{]}.
A preprint audit of sixteen deployed agent frameworks reports none that
writes a run record a reader can check without trusting the writer; a
language model scored the audit grid, and agreement with human coders
fell below the registered threshold. It proposes a log kept outside the
agent's harness {[}\protect\hyperlink{ref-Dai2026Hearsay}{40}{]}. A
coverage model separates event time from verifier receipt time, judges
whether a verifier's evidence was complete at a stated cutoff, links
records to decisions and distinguishes refused execution branches from
late reports; it also separates absence from one retrieval from absence
in a verifier view whose closure, including retrieval coverage and
retention, is justified, and specifies that basis without designing its
implementation {[}\protect\hyperlink{ref-Wu2026EvidenceCoverage}{31}{]}.
& The gateway outside the agent attaches a coverage record to every
result: the entities and interval requested and evaluated, the as-of
cutoff, each source's event and ingestion watermarks and known gaps, the
retention horizon, limits reached, access filtering, sampling and schema
versions. Fixed rules outside the agent derive whether a result is
complete, partial, truncated, denied or unavailable, and a scope-limited
flag confines absence claims to the authorized scope (manuscript Section
6.2). Every fact, including a derived record, keeps event time and the
time it first became available to queries (S4.1). \\
C3: human release and dependence & Truth-maintenance systems record each
belief's justifications and retract conclusions whose support changes
{[}\protect\hyperlink{ref-Doyle1979TMS}{41},
\protect\hyperlink{ref-Malheiro1994BeliefRevision}{42}{]}; STIX keeps
observed data apart from intelligence assertions and opinions
{[}\protect\hyperlink{ref-OASIS2021STIX}{43}{]}, and US intelligence
analytic standards ask analytic products to distinguish underlying
information from assumptions and judgments
{[}\protect\hyperlink{ref-ODNI2015ICD203}{44}{]}. Anamnesis starts
model-derived claims as candidates, and no memory operation can raise
their authority without a separately recorded grant; the grantor need
not be a human {[}\protect\hyperlink{ref-Dan2026Anamnesis}{38}{]}.
GovMem groups dependent agent traces when estimating support instead of
counting them as independent votes, permits automatic promotion and
routes unresolved candidates to human review
{[}\protect\hyperlink{ref-Qi2026GovMem}{45}{]}. A data-governance
preprint admits only verified records as precedent, without saying who
authorizes promotion or what a precedent may establish (S9.4)
{[}\protect\hyperlink{ref-Akhtar2026Contestable}{46}{]}. Albarrak and
Jagtap place human approval before intelligence is shared, but do not
state that analyses kept in their institutional-memory store pass that
gate {[}\protect\hyperlink{ref-AlbarrakJagtap2026Industry4CTI}{47}{]}.
MAP-Graph writes memories without human release (S9.3)
{[}\protect\hyperlink{ref-Wang2026MAPGraph}{48}{]}; Li and coauthors
check an authorization state that need not come from a person
{[}\protect\hyperlink{ref-Li2026EvidenceToEffect}{49}{]}. A formal study
distinguishes independent observations from reports derived from shared
evidence and requires a dependence model as well as ancestry
{[}\protect\hyperlink{ref-Bara2026EpistemicSybil}{50}{]}; a
threat-detection design drops its own compiled detections at ingest
(S9.4) {[}\protect\hyperlink{ref-Evron2026AdversaryBrains}{51}{]}; an
engineering assurance design binds evidence to the exact artifact, and
its governing profile can require human authority
{[}\protect\hyperlink{ref-Koch2026AgileV}{52}{]}. AgenticCyOps provides
analyst oversight without a human release decision for every item before
organizational reuse
{[}\protect\hyperlink{ref-Mitra2026AgenticCyOps}{53}{]}. TOKI's optional
external callback, which may be a human, only chooses between
contradicting facts {[}\protect\hyperlink{ref-Wang2026TOKI}{36}{]}. &
The recorded grant is specialized to a human principal for every
conclusion entering organizational memory, as the data-governance gate
does for precedent, and released items stay context: a released
conclusion is never counted as a new or independent observation, and a
claim about current conduct must cite observations of that conduct from
the current period, not scores or control decisions alone (manuscript
Sections 4.2 and 6.4). Labels, assessments, control decisions,
detections and threat-knowledge items carry a link to every agent output
and every released organizational-memory item open in the writer's case
or session. The labeling, assessment, control, detection and
threat-knowledge interfaces, rather than the person, write these links
automatically, so they also cover records a person writes with an agent
output open in the case. Linked records count as dependent; a missing
link does not establish independence (manuscript Section 4.2).
Responsibilities 21 and 22, quality review and human release, are
inherited from threat-intelligence work (S4). \\

\bottomrule\noalign{}
\end{tabular}
}
\end{table}

\hypertarget{s9.1-payment-decisions-and-the-outcomes-they-reveal}{%
\subsubsection{S9.1 Payment decisions and the outcomes they
reveal}\label{s9.1-payment-decisions-and-the-outcomes-they-reveal}}

Payment engineering already connects records, investigation and
protective decisions. Stripe describes similarity clustering for
investigators {[}\protect\hyperlink{ref-Tausz2020}{54}{]}. It also
describes a retraining cycle for card-testing models: labels are derived
partly through manual expert review, and retrained models are evaluated
on offline data and compared with the old model in blue-green tests
before deployment {[}\protect\hyperlink{ref-Drapeau2024}{55}{]}. These
vendor accounts establish publicly described practices in their stated
settings. They do not independently measure the effectiveness of the
proposed specification.

Outcome-aware evaluation also predates this work. Manapat's 2015 talk
abstract describes adding randomness to production fraud scoring so that
a blocking model's precision and recall can be estimated from production
data {[}\protect\hyperlink{ref-Manapat2015Counterfactual}{32}{]}. Adyen
evaluates suggested payment interventions using transactions logged
under an existing policy
{[}\protect\hyperlink{ref-Egg2025AdyenOPE}{27}{]}. A Meta preprint
proposes estimators for blocking's false-positive rate, which cannot be
monitored directly after a user is blocked, and tests them on simulated
blocking policies
{[}\protect\hyperlink{ref-Tocker2022VirtualControl}{28}{]}. These
approaches motivate retaining the decision and the conditions under
which its outcome became observable. They do not make an unreported
outcome known.

Two recent proposals connect evaluation evidence to operational
authority. Deng treats fraud automation as an authorization decision
because its evaluation labels are selective and delayed
{[}\protect\hyperlink{ref-Deng2026FCAC}{29}{]}. Solozobov specifies
decision-linked records for facts learned later, including the arrival
of ground truth; the schema seals each decision record and types the
links between decisions
{[}\protect\hyperlink{ref-Solozobov2026DecisionTrace}{30}{]}. C1 applies
that linkage to several outcome channels. It keeps the decision as the
unit of record and adds each outcome report as a separate appended
record, linked to every decision that may have shaped it with that
decision's role. It records what a control closed, delayed or changed in
each channel, alongside reports and assessments of their truth. Its
contribution is this payment-investigation record contract, not a new
estimator of missing outcomes.

\hypertarget{s9.2-intelligence-production-and-investigation}{%
\subsubsection{S9.2 Intelligence production and
investigation}\label{s9.2-intelligence-production-and-investigation}}

STIX is a structured language for representing and exchanging cyber
threat intelligence {[}\protect\hyperlink{ref-OASIS2021STIX}{43}{]}.
Such representations make analysis reusable. STIX already separates
observed data from intelligence assertions and opinions, but it does not
address how opinions should be interpreted
{[}\protect\hyperlink{ref-OASIS2021STIX}{43}{]}. C3 supplies a rule for
what a reused conclusion can establish in this specification.

ExCyTIn-Bench evaluates agents on multi-step questions over security
logs {[}\protect\hyperlink{ref-Wu2026ExCyTIn}{56}{]}. This benchmark
establishes relevant investigative tasks and a way to evaluate
particular capabilities. The present specification adds requirements for
interpreting the inputs to such tasks and preserving the status of their
outputs across payment products.

PROVSEEK is especially close to the requirement that a claim remain
checkable. Its agent cites audit-event identifiers for claims and
reports searched-for artifacts it did not find
{[}\protect\hyperlink{ref-Mukherjee2025PROVSEEK}{57}{]}. C2 additionally
requires a gateway outside the agent to record query coverage. A missing
artifact then has an explicit search scope, rather than relying on the
agent's description of what it searched.

\hypertarget{s9.3-permission-to-act-and-authority-to-reuse}{%
\subsubsection{S9.3 Permission to act and authority to
reuse}\label{s9.3-permission-to-act-and-authority-to-reuse}}

Recent proposals bind authority to particular artifacts or actions.
MAP-Graph represents memories, claims and actions in a provenance graph
and records each action-gate decision with its supporting memories; the
gate can return an AskUser decision for some high-risk actions
{[}\protect\hyperlink{ref-Wang2026MAPGraph}{48}{]}. Li and colleagues
relate an execution attempt and its artifact to a downstream-use
authorization state, checked at publication
{[}\protect\hyperlink{ref-Li2026EvidenceToEffect}{49}{]}. Agile-V binds
evidence to an exact artifact and links effects to authorization
records; its governing profile can require human authority
{[}\protect\hyperlink{ref-Koch2026AgileV}{52}{]}. These are concrete
precedents for authorization records. C3 makes a narrower, stronger
choice for organizational memory: a person releases every conclusion,
and the released conclusion remains context or hypothesis.

A survey of persistent agents calls for records tying actions to the
state that justified them
{[}\protect\hyperlink{ref-Ding2026AlwaysOn}{58}{]}. EDDOps links
versioned evaluation records to the agent changes they prompt
{[}\protect\hyperlink{ref-Xia2024EDDOps}{59}{]}. Both motivate
preserving the connection between an output and its later effects. Under
C3, labels, assessments, controls, detections and threat-knowledge items
link back to every agent output and every released organizational-memory
item open in the writer's case or session. Human release does not erase
this dependence or turn the output into independent evidence.

\hypertarget{s9.4-memory-dependence-and-revision}{%
\subsubsection{S9.4 Memory, dependence and
revision}\label{s9.4-memory-dependence-and-revision}}

Agents can learn reusable lessons from experience. ExpeL extracts
lessons, which it calls insights, from its successes and failures on
training tasks, adds them to the prompt for later tasks and retrieves
similar successful trajectories as examples
{[}\protect\hyperlink{ref-Zhao2024ExpeL}{60}{]}. This capability
explains the value of reuse, but reuse can propagate an earlier mistake.
A-MemGuard checks memory-derived reasoning during use, consulting
lessons from stored failures
{[}\protect\hyperlink{ref-Wei2026AMemGuard}{61}{]}. Manuscript Section 7
proposes testing new lessons against counterexamples before reuse, then
testing their effects in E7. This proposed control has no measured
payment-platform benefit here.

The choice of memory structure depends on purpose. An investigative
system needs to know whether a retrieved item is an observation or a
conclusion. Studies of memory consolidation show that it can erase a
claim's source; external source records restore much of the lost control
in their test settings
{[}\protect\hyperlink{ref-Zhan2026AuthMemBench}{62},
\protect\hyperlink{ref-Xu2026PPMF}{63}{]}. Those settings used synthetic
customer-service histories and benchmark attacks, not payment
investigations.

SSGM separates an append-only observation log from derived memory and
checks updates against protected facts
{[}\protect\hyperlink{ref-Lam2026SSGM}{64}{]}. A shared-memory service
records writers, derivation links and supersession; its authors evaluate
it without a baseline
{[}\protect\hyperlink{ref-Margalit2026GovernedMemory}{65}{]}. PROV-AGENT
extends provenance records to prompts, responses and decisions, linking
them to downstream outcomes
{[}\protect\hyperlink{ref-Souza2025ProvAgent}{66}{]}. These precedents
support typed records and traceable revision. C3 does not claim to
invent them.

Source ancestry is necessary but not sufficient to establish independent
support. A formal treatment of evidence dependence distinguishes
independent observations from reports derived from shared evidence
{[}\protect\hyperlink{ref-Bara2026EpistemicSybil}{50}{]}. C3 follows the
same concern: several agents repeating one report do not create several
independent observations. A released conclusion can guide the next
query, but evidence for that conclusion must remain traceable to the
underlying records.

Two recent designs bear directly on C3. A data-governance preprint
commits no agent-proposed metadata decision before a steward reviews it.
It proposes a promotion gate that admits only verified records as
precedent, records each promotion and lets a reviewer demote a record,
flagging the decisions that relied on it. It does not say who authorizes
promotion or what a retrieved precedent may establish
{[}\protect\hyperlink{ref-Akhtar2026Contestable}{46}{]}. A
threat-detection design tags every detection compiled from one of its
hypotheses and drops tagged observations at ingest, so its own
conclusions cannot re-enter as fresh evidence. It leaves as a residual
risk an analyst who writes such an alert into a report that is later
ingested, and its reference implementation does not yet exercise the
rule {[}\protect\hyperlink{ref-Evron2026AdversaryBrains}{51}{]}. C3
keeps released conclusions as context rather than evidence, and its
interface-written links cover records a person writes with an agent
output open in the case.

\hypertarget{s9.5-time-and-coverage}{%
\subsubsection{S9.5 Time and coverage}\label{s9.5-time-and-coverage}}

Timing matters when it changes which information a decision could use.
Dataflow addresses out-of-order streams, where the system cannot assume
that all events have arrived
{[}\protect\hyperlink{ref-Akidau2015Dataflow}{35}{]}. Leakage research
identifies information about a target that should not legitimately be
available to a model
{[}\protect\hyperlink{ref-Kaufman2011Leakage}{67}{]}. C2 uses these
concerns to distinguish when a fact held from when it became queryable.
A relationship learned tomorrow cannot support a decision frozen today
merely because the relationship existed yesterday.

Recent designs combine timing with governance. Quipu controls graph
writes while retaining validity and recording times; some policies
require human approval bound to the evidence
{[}\protect\hyperlink{ref-Brown2026Quipu}{39}{]}. Wu's coverage model
distinguishes event time from verifier receipt time, judges whether a
verifier's evidence was complete at a stated cutoff, links records to
decisions, and separates refused execution branches from late reports
{[}\protect\hyperlink{ref-Wu2026EvidenceCoverage}{31}{]}. It also
separates absence from one retrieval from absence in a verifier view
whose closure, including retrieval coverage and retention, is justified,
and it specifies that basis without designing its implementation. C2
applies these ideas to a query contract that records scope, coverage and
availability outside the investigating agent.

Together, these precedents explain why C1 to C3 are a synthesis. The
proposed contribution is their joint application to payment-platform
investigation, with explicit channels, record classes, query responses
and release rules. The hypothetical traces expose conclusions that fail
when specified context is removed. E1 to E9 test parts of that context,
including whether the channel states reduce bias relative to a
platform's current label rule (E9) and, through planted drafts that rest
only on released items, the rule that a released conclusion adds no
independent corroboration (E8; manuscript Section 8.3). In S7.1, E5 and
E8 measure reviewer time, and no test measures the cost of collecting
the records (manuscript Section 9.1). Neither the comparison nor the
traces establish that the integrated system works in deployment.

\FloatBarrier

\hypertarget{references}{%
\subsection{References}\label{references}}

\suppressfloats[t]

\hypertarget{refs}{}
\begin{CSLReferences}{1}{1}
\leavevmode\vadjust pre{\hypertarget{ref-MITRE2026F3}{}}%
{[}1{]} MITRE Center for Threat-Informed Defense. 2026. Fight Fraud
Framework (F3), version 1.1. GitHub,
center-for-threat-informed-defense/fight-fraud-framework. Cited version:
commit cb186c2 (18 September 2026).
\url{https://github.com/center-for-threat-informed-defense/fight-fraud-framework/tree/cb186c2b3372164e0b7def0d4f09531b6acc16e5}.
Accessed 25 October 2026.

\leavevmode\vadjust pre{\hypertarget{ref-Meli2019}{}}%
{[}2{]} M. Meli, M. R. McNiece, and B. Reaves. 2019. How Bad Can It Git?
Characterizing Secret Leakage in Public GitHub Repositories.
\emph{Network and Distributed Systems Security (NDSS) Symposium 2019}.
\url{https://doi.org/10.14722/ndss.2019.23418}.

\leavevmode\vadjust pre{\hypertarget{ref-Demir2026}{}}%
{[}3{]} N. Demir, Y. Vekaria, G. Smaragdakis, and Z. Durumeric. 2026.
Keys on Doormats: Exposed API Credentials on the Web. arXiv preprint
2603.12498v3. Accepted at ACM CCS 2026 (November 2026) per the reference
printed in arXiv v3. \url{https://arxiv.org/abs/2603.12498v3}.

\leavevmode\vadjust pre{\hypertarget{ref-Gamazo2023}{}}%
{[}4{]} W. Gamazo and N. Quist. 2023. CloudKeys in the Air: Tracking
Malicious Operations of Exposed IAM Keys. Palo Alto Networks Unit 42.
\url{https://unit42.paloaltonetworks.com/malicious-operations-of-exposed-iam-keys-cryptojacking/}.
Accessed 19 October 2026.

\leavevmode\vadjust pre{\hypertarget{ref-Microsoft2022AiTM}{}}%
{[}5{]} Microsoft Threat Intelligence. 2022. From cookie theft to BEC:
Attackers use AiTM phishing sites as entry point to further financial
fraud. Microsoft Security Blog.
\url{https://www.microsoft.com/en-us/security/blog/2022/07/12/from-cookie-theft-to-bec-attackers-use-aitm-phishing-sites-as-entry-point-to-further-financial-fraud/}.
Accessed 19 October 2026.

\leavevmode\vadjust pre{\hypertarget{ref-Stripe2026MonitoringPrograms}{}}%
{[}6{]} Stripe. 2026. Dispute and fraud card monitoring programs. Stripe
documentation.
\url{https://docs.stripe.com/disputes/monitoring-programs}. Accessed 25
October 2026.

\leavevmode\vadjust pre{\hypertarget{ref-Stripe2026Disputes}{}}%
{[}7{]} Stripe. 2026. How disputes work. Stripe documentation.
\url{https://docs.stripe.com/disputes/how-disputes-work}. Accessed 25
October 2026.

\leavevmode\vadjust pre{\hypertarget{ref-Stripe2026SEPADebit}{}}%
{[}8{]} Stripe. 2026. SEPA Direct Debit payments. Stripe documentation.
\url{https://docs.stripe.com/payments/sepa-debit}. Accessed 25 October
2026.

\leavevmode\vadjust pre{\hypertarget{ref-Stripe2026ACHTerms}{}}%
{[}9{]} Stripe. 2026. ACH Payment Terms. Stripe Legal.
\url{https://stripe.com/legal/ACH}. Accessed 25 October 2026.

\leavevmode\vadjust pre{\hypertarget{ref-LeadBank2026ACHReturns}{}}%
{[}10{]} Lead Bank. 2026. Returns, Reversals, \& Exceptions. Lead Bank
documentation.
\url{https://docs.lead.bank/api-integrations/ach/return-codes}. Accessed
25 October 2026.

\leavevmode\vadjust pre{\hypertarget{ref-StripeOpenAPI2026}{}}%
{[}11{]} Stripe. 2026. OpenAPI specification for the Stripe API. GitHub,
stripe/openapi.
\url{https://github.com/stripe/openapi/tree/18fa2cc768024b47789dc6fa243acd48478acac2}.
Accessed 22 October 2026.

\leavevmode\vadjust pre{\hypertarget{ref-Leach2019}{}}%
{[}12{]} B. Leach. 2019. Fast and flexible observability with canonical
log lines. Stripe blog.
\url{https://stripe.com/blog/canonical-log-lines}. Accessed 21 October
2026.

\leavevmode\vadjust pre{\hypertarget{ref-Dhama2026Causal}{}}%
{[}13{]} G. Dhama. 2026. Causal Label Recovery in Payment Networks.
arXiv preprint 2605.29272v1. \url{https://arxiv.org/abs/2605.29272v1}.

\leavevmode\vadjust pre{\hypertarget{ref-Liu2018GEM}{}}%
{[}14{]} Z. Liu, C. Chen, X. Yang, J. Zhou, X. Li, and L. Song. 2018.
Heterogeneous Graph Neural Networks for Malicious Account Detection.
\emph{Proceedings of the 27th ACM International Conference on
Information and Knowledge Management}.
\url{https://doi.org/10.1145/3269206.3272010}.

\leavevmode\vadjust pre{\hypertarget{ref-Fatima2024}{}}%
{[}15{]} K. Fatima, M. Schuckers, G. Cruz-Ortiz, D. Hou, S. Purnapatra,
T. Andrews, A. Neupane, B. Marshall, and S. Schuckers. 2024. A
Large-Scale Study of Performance and Equity of Commercial Remote
Identity Verification Technologies Across Demographics. \emph{2024 IEEE
International Joint Conference on Biometrics}, pp. 1--8.
\url{https://doi.org/10.1109/IJCB62174.2024.10744432}.

\leavevmode\vadjust pre{\hypertarget{ref-Doerfler2019}{}}%
{[}16{]} P. Doerfler, M. Marincenko, J. Ranieri, Y. Jiang, A. Moscicki,
D. McCoy, and K. Thomas. 2019. Evaluating Login Challenges as a Defense
Against Account Takeover. \emph{Proceedings of the 2019 World Wide Web
Conference}, pp. 372--382.
\url{https://doi.org/10.1145/3308558.3313481}.

\leavevmode\vadjust pre{\hypertarget{ref-Sharma2026AuditableFraud}{}}%
{[}17{]} R. Sharma. 2026. Toward Auditable Fraud Detection: Combining
Graph Features, Model Explanations, and Agentic Case Investigation.
arXiv preprint 2607.19266v1. \url{https://arxiv.org/abs/2607.19266v1}.

\leavevmode\vadjust pre{\hypertarget{ref-Nasr2025}{}}%
{[}18{]} M. Nasr, N. Carlini, C. Sitawarin, S. V. Schulhoff, J. Hayes,
M. Ilie, J. Pluto, S. Song, H. Chaudhari, I. Shumailov, A. Thakurta, K.
Y. Xiao, A. Terzis, and F. Tramèr. 2025. The Attacker Moves Second:
Stronger Adaptive Attacks Bypass Defenses Against LLM Jailbreaks and
Prompt Injections. arXiv preprint 2510.09023v1.
\url{https://arxiv.org/abs/2510.09023v1}.

\leavevmode\vadjust pre{\hypertarget{ref-Debenedetti2025CaMeL}{}}%
{[}19{]} E. Debenedetti, I. Shumailov, T. Fan, J. Hayes, N. Carlini, D.
Fabian, C. Kern, C. Shi, A. Terzis, and F. Tramèr. 2025. Defeating
Prompt Injections by Design. arXiv preprint 2503.18813v2.
\url{https://arxiv.org/abs/2503.18813v2}.

\leavevmode\vadjust pre{\hypertarget{ref-Dash2026}{}}%
{[}20{]} P. Dash, T. Ge, A. Jain, T. Shah, and Z. Shang. 2026. From
Untrusted Input to Trusted Memory: A Systematic Study of Memory
Poisoning Attacks in LLM Agents. arXiv preprint 2606.04329v2.
\url{https://arxiv.org/abs/2606.04329v2}.

\leavevmode\vadjust pre{\hypertarget{ref-Xiong2026MemoryManagement}{}}%
{[}21{]} Z. Xiong, Y. Lin, W. Xie, P. He, Z. Liu, J. Tang, H. Lakkaraju,
and Z. Xiang. 2026. How Memory Management Impacts LLM Agents: An
Empirical Study of Experience-Following Behavior. \emph{Proceedings of
the 64th Annual Meeting of the Association for Computational Linguistics
(Volume 1: Long Papers)}, pp. 623--645.
\url{https://doi.org/10.18653/v1/2026.acl-long.27}.

\leavevmode\vadjust pre{\hypertarget{ref-Lakkaraju2017Selective}{}}%
{[}22{]} H. Lakkaraju, J. Kleinberg, J. Leskovec, J. Ludwig, and S.
Mullainathan. 2017. The Selective Labels Problem: Evaluating Algorithmic
Predictions in the Presence of Unobservables. \emph{Proceedings of the
23rd ACM SIGKDD International Conference on Knowledge Discovery and Data
Mining}, pp. 275--284. \url{https://doi.org/10.1145/3097983.3098066}.

\leavevmode\vadjust pre{\hypertarget{ref-Bottou2013Counterfactual}{}}%
{[}23{]} L. Bottou, J. Peters, J. Quiñonero-Candela, D. X. Charles, D.
M. Chickering, E. Portugaly, D. Ray, P. Simard, and E. Snelson. 2013.
Counterfactual Reasoning and Learning Systems: The Example of
Computational Advertising. \emph{Journal of Machine Learning Research},
14(101), pp. 3207--3260.
\url{https://jmlr.org/papers/v14/bottou13a.html}.

\leavevmode\vadjust pre{\hypertarget{ref-Singh2019DecisionProvenance}{}}%
{[}24{]} J. Singh, J. Cobbe, and C. Norval. 2019. Decision Provenance:
Harnessing Data Flow for Accountable Systems. \emph{IEEE Access}, 7, pp.
6562--6574. \url{https://doi.org/10.1109/ACCESS.2018.2887201}.

\leavevmode\vadjust pre{\hypertarget{ref-Ehrhardt2021RejectInference}{}}%
{[}25{]} A. Ehrhardt, C. Biernacki, V. Vandewalle, P. Heinrich, and S.
Beben. 2021. Reject Inference Methods in Credit Scoring. \emph{Journal
of Applied Statistics}, 48(13-15), pp. 2734--2754.
\url{https://doi.org/10.1080/02664763.2021.1929090}.

\leavevmode\vadjust pre{\hypertarget{ref-Dhama2026Limits}{}}%
{[}26{]} G. Dhama. 2026. The Fundamental Limits of Fraud Detection in
Card Payment Networks: An Information-Theoretic Theory of Delayed,
Censored, Corrupted, and Counterfactual Feedback. arXiv preprint
2605.27557v2. \url{https://arxiv.org/abs/2605.27557v2}.

\leavevmode\vadjust pre{\hypertarget{ref-Egg2025AdyenOPE}{}}%
{[}27{]} A. Egg. 2025. Off-policy Evaluation for Payments at Adyen.
arXiv preprint 2501.10470v1. \url{https://arxiv.org/abs/2501.10470v1}.

\leavevmode\vadjust pre{\hypertarget{ref-Tocker2022VirtualControl}{}}%
{[}28{]} M. Tocker. 2022. Virtual Control Group: Measuring Hidden
Performance Metrics. arXiv preprint 2208.12941v1.
\url{https://arxiv.org/abs/2208.12941v1}.

\leavevmode\vadjust pre{\hypertarget{ref-Deng2026FCAC}{}}%
{[}29{]} J. Deng. 2026. When Can Fraud Operations Authorize Automation?
A Decision-Support Framework for Fresh Audit Evidence and Review
Workload. arXiv preprint 2608.08577v1.
\url{https://arxiv.org/abs/2608.08577v1}.

\leavevmode\vadjust pre{\hypertarget{ref-Solozobov2026DecisionTrace}{}}%
{[}30{]} O. Solozobov. 2026. Decision Trace Schema for Governance
Evidence in Real-Time Risk Systems. arXiv preprint 2604.09296v1.
\url{https://arxiv.org/abs/2604.09296v1}.

\leavevmode\vadjust pre{\hypertarget{ref-Wu2026EvidenceCoverage}{}}%
{[}31{]} M. Wu, L. Wang, Y. Zhang, and J. Deng. 2026. Evidence Coverage
for Intent-Bound Execution: Scope, Obligations, and Cutoff Reasoning.
arXiv preprint 2610.01386v1. \url{https://arxiv.org/abs/2610.01386v1}.

\leavevmode\vadjust pre{\hypertarget{ref-Manapat2015Counterfactual}{}}%
{[}32{]} M. Manapat. 2015. Counterfactual evaluation of machine learning
models. Conference talk, PyData Seattle.
\url{https://pydata.org/seattle2015/schedule/presentation/17/}.

\leavevmode\vadjust pre{\hypertarget{ref-Snodgrass1986Temporal}{}}%
{[}33{]} R. Snodgrass and I. Ahn. 1986. Temporal Databases.
\emph{Computer}, 19(9), pp. 35--42.
\url{https://doi.org/10.1109/MC.1986.1663327}.

\leavevmode\vadjust pre{\hypertarget{ref-Razniewski2011Completeness}{}}%
{[}34{]} S. Razniewski and W. Nutt. 2011. Completeness of Queries over
Incomplete Databases. \emph{Proceedings of the VLDB Endowment}, 4(11),
pp. 749--760. \url{https://doi.org/10.14778/3402707.3402715}.

\leavevmode\vadjust pre{\hypertarget{ref-Akidau2015Dataflow}{}}%
{[}35{]} T. Akidau, R. Bradshaw, C. Chambers, S. Chernyak, R. J.
Fernández-Moctezuma, R. Lax, S. McVeety, D. Mills, F. Perry, E. Schmidt,
and S. Whittle. 2015. The Dataflow Model: A Practical Approach to
Balancing Correctness, Latency, and Cost in Massive-Scale, Unbounded,
Out-of-Order Data Processing. \emph{Proceedings of the VLDB Endowment},
8(12), pp. 1792--1803. \url{https://doi.org/10.14778/2824032.2824076}.

\leavevmode\vadjust pre{\hypertarget{ref-Wang2026TOKI}{}}%
{[}36{]} Z. Wang. 2026. TOKI: A Bitemporal Operator Algebra for
Contradiction Resolution in LLM-Agent Persistent Memory. arXiv preprint
2606.06240v1. \url{https://arxiv.org/abs/2606.06240v1}.

\leavevmode\vadjust pre{\hypertarget{ref-Rasmussen2025Zep}{}}%
{[}37{]} P. Rasmussen, P. Paliychuk, T. Beauvais, J. Ryan, and D.
Chalef. 2025. Zep: A Temporal Knowledge Graph Architecture for Agent
Memory. arXiv preprint 2501.13956v1.
\url{https://arxiv.org/abs/2501.13956v1}.

\leavevmode\vadjust pre{\hypertarget{ref-Dan2026Anamnesis}{}}%
{[}38{]} A. Dan. 2026. Anamnesis: An Authority-Preserving Runtime
Contract for Persistent AI Memory. Zenodo preprint, Version 1.0, 16
August 2026. \url{https://zenodo.org/records/22118090}.

\leavevmode\vadjust pre{\hypertarget{ref-Brown2026Quipu}{}}%
{[}39{]} S. Brown. 2026. Quipu: A Governed Bitemporal Knowledge Graph
Store. arXiv preprint 2608.16813v1.
\url{https://arxiv.org/abs/2608.16813v1}.

\leavevmode\vadjust pre{\hypertarget{ref-Dai2026Hearsay}{}}%
{[}40{]} J. Dai, Z. Yang, P. Shao, K. Ng, Z. Liu, C. Ju, Y. He, and B.
Hu. 2026. Hearsay: Can an Auditor Trust the Record a Deployed Agent
Harness Writes? arXiv preprint 2609.32495v1.
\url{https://arxiv.org/abs/2609.32495v1}.

\leavevmode\vadjust pre{\hypertarget{ref-Doyle1979TMS}{}}%
{[}41{]} J. Doyle. 1979. A Truth Maintenance System. \emph{Artificial
Intelligence}, 12(3), pp. 231--272.
\url{https://doi.org/10.1016/0004-3702(79)90008-0}.

\leavevmode\vadjust pre{\hypertarget{ref-Malheiro1994BeliefRevision}{}}%
{[}42{]} B. Malheiro, N. R. Jennings, and E. Oliveira. 1994. Belief
Revision in Multi-Agent Systems. \emph{Proceedings of the 11th European
Conference on Artificial Intelligence}, pp. 294--298.
\url{https://eprints.soton.ac.uk/252143}.

\leavevmode\vadjust pre{\hypertarget{ref-OASIS2021STIX}{}}%
{[}43{]} OASIS. 2021. STIX Version 2.1. OASIS Standard.
\url{https://docs.oasis-open.org/cti/stix/v2.1/os/stix-v2.1-os.html}.

\leavevmode\vadjust pre{\hypertarget{ref-ODNI2015ICD203}{}}%
{[}44{]} Office of the Director of National Intelligence. 2015.
Intelligence Community Directive 203: Analytic Standards. Directive. As
amended by technical amendments through 21 January 2022.
\url{https://archive.dni.gov/files/documents/ICD/ICD-203.pdf}.

\leavevmode\vadjust pre{\hypertarget{ref-Qi2026GovMem}{}}%
{[}45{]} Y. Qi, X. Xu, and Y. Li. 2026. When Not to Write Memory:
Governing False Promotion from Correlated Agent Traces. arXiv preprint
2607.02579v1. \url{https://arxiv.org/abs/2607.02579v1}.

\leavevmode\vadjust pre{\hypertarget{ref-Akhtar2026Contestable}{}}%
{[}46{]} H. Akhtar and B. Freeman. 2026. Contestable Agents for Data
Governance: Recording Human Stewardship of AI-Generated Metadata Across
Structured and Unstructured Systems. arXiv preprint 2610.31553v1.
\url{https://arxiv.org/abs/2610.31553v1}.

\leavevmode\vadjust pre{\hypertarget{ref-AlbarrakJagtap2026Industry4CTI}{}}%
{[}47{]} M. Albarrak and S. Jagtap. 2026. Designing an AI-Assisted Cyber
Threat Intelligence Framework for Industry 4.0: A Human-in-the-Loop
Design Science Approach. \emph{Applied Sciences}, 16(15), Article 7646.
\url{https://doi.org/10.3390/app16157646}.

\leavevmode\vadjust pre{\hypertarget{ref-Wang2026MAPGraph}{}}%
{[}48{]} Y. Wang, Z. Yan, J. Zhang, Z. Wu, M. Zheng, Z. Sun, Y. Zhu, and
T. Cai. 2026. MAP-Graph: Provenance-Aware Shared Memory for Multi-Agent
Workflows. arXiv preprint 2608.10509v1.
\url{https://arxiv.org/abs/2608.10509v1}.

\leavevmode\vadjust pre{\hypertarget{ref-Li2026EvidenceToEffect}{}}%
{[}49{]} Y. Li, Z. Xu, S. Volkov, H. Liu, T. Zhou, X. Chen, D. Wan, D.
Shao, Y. Luo, and H. Sun. 2026. From Evidence to Effect: Authority
Semantics and Runtime Infrastructure for Stateful Agents. arXiv preprint
2609.08472v3. Cited version: arXiv 2609.08472v3 (28 September 2026).
\url{https://arxiv.org/abs/2609.08472v3}.

\leavevmode\vadjust pre{\hypertarget{ref-Bara2026EpistemicSybil}{}}%
{[}50{]} M. Bara. 2026. Epistemic Sybil Resistance: Multiplying AI
Agents Without Multiplying Evidence. arXiv preprint 2609.01873v1.
\url{https://arxiv.org/abs/2609.01873v1}.

\leavevmode\vadjust pre{\hypertarget{ref-Evron2026AdversaryBrains}{}}%
{[}51{]} G. Evron. 2026. Adversary Brains: Hypothesis-Based Logical
Signatures for Threat Actor Detection with LLMs. arXiv preprint
2610.35624v1. \url{https://arxiv.org/abs/2610.35624v1}.

\leavevmode\vadjust pre{\hypertarget{ref-Koch2026AgileV}{}}%
{[}52{]} C. Koch. 2026. From Agent Output to Authorized Transition.
arXiv preprint 2609.28216v1. The v1 PDF title page reads ``From Agent
Output to Authorized Transition: The Agile-V Assurance Spine for
State-Bound Evidence in Software, Firmware, and Hardware Engineering''
instead. \url{https://arxiv.org/abs/2609.28216v1}.

\leavevmode\vadjust pre{\hypertarget{ref-Mitra2026AgenticCyOps}{}}%
{[}53{]} S. Mitra, R. Patel, S. Mittal, M. R. Rahman, and S. Rahimi.
2026. AgenticCyOps: Securing Multi-Agentic AI Integration in Enterprise
Cyber Operations. arXiv preprint 2603.09134v1.
\url{https://arxiv.org/abs/2603.09134v1}.

\leavevmode\vadjust pre{\hypertarget{ref-Tausz2020}{}}%
{[}54{]} A. Tausz. 2020. Similarity clustering to catch fraud rings.
Stripe blog. \url{https://stripe.com/blog/similarity-clustering}.
Accessed 21 October 2026.

\leavevmode\vadjust pre{\hypertarget{ref-Drapeau2024}{}}%
{[}55{]} R. Drapeau and I. Khosla. 2024. The ML flywheel: how we
continually improve our models to reduce card testing. Stripe blog.
\url{https://stripe.com/blog/the-ml-flywheel-how-we-continually-improve-our-models-to-reduce-card-testing}.
Accessed 21 October 2026.

\leavevmode\vadjust pre{\hypertarget{ref-Wu2026ExCyTIn}{}}%
{[}56{]} Y. Wu, M. Velazco, A. Zhao, M. R. Meléndez Luján, S. Movva, Y.
K. Roy, Q. Nguyen, R. Rodriguez, Q. Wu, M. Albada, J. Kiseleva, and A.
Mudgerikar. 2026. ExCyTIn-Bench: Evaluating LLM agents on Cyber Threat
Investigation. arXiv preprint 2507.14201v3. Accepted at ICML 2026; final
proceedings pagination not established. Cited version: arXiv
2507.14201v3 (1 May 2026). \url{https://arxiv.org/abs/2507.14201v3}.

\leavevmode\vadjust pre{\hypertarget{ref-Mukherjee2025PROVSEEK}{}}%
{[}57{]} K. Mukherjee and M. Kantarcioglu. 2025. LLM-driven Provenance
Forensics for Threat Investigation and Detection. arXiv preprint
2508.21323v2. The v2 PDF title page reads ``LLM-driven Provenance
Forensics for Threat Intelligence and Detection'' instead.
\url{https://arxiv.org/abs/2508.21323v2}.

\leavevmode\vadjust pre{\hypertarget{ref-Ding2026AlwaysOn}{}}%
{[}58{]} T. Ding, A. Nannapaneni, B. Liu, and L. Zhang. 2026. Always-On
Agents: A Survey of Persistent Memory, State, and Governance in LLM
Agents. arXiv preprint 2606.30306v1.
\url{https://arxiv.org/abs/2606.30306v1}.

\leavevmode\vadjust pre{\hypertarget{ref-Xia2024EDDOps}{}}%
{[}59{]} B. Xia, Q. Lu, L. Zhu, Z. Xing, D. Zhao, and H. Zhang. 2025.
Evaluation-Driven Development and Operations of LLM Agents: A Process
Model and Reference Architecture. arXiv preprint 2411.13768v3.
\url{https://arxiv.org/abs/2411.13768v3}.

\leavevmode\vadjust pre{\hypertarget{ref-Zhao2024ExpeL}{}}%
{[}60{]} A. Zhao, D. Huang, Q. Xu, M. Lin, Y.-J. Liu, and G. Huang.
2024. ExpeL: LLM Agents Are Experiential Learners. \emph{Proceedings of
the AAAI Conference on Artificial Intelligence}, 38(17), pp.
19632--19642. \url{https://doi.org/10.1609/aaai.v38i17.29936}.

\leavevmode\vadjust pre{\hypertarget{ref-Wei2026AMemGuard}{}}%
{[}61{]} Q. Wei, T. Yang, Y. Wang, X. Li, L. Li, Z. Yin, Y. Zhan, T.
Holz, Z. Lin, and X. Wang. 2025. A-MemGuard: A Proactive Defense
Framework for LLM-Based Agent Memory. arXiv preprint 2510.02373v1. ICML
2026 acceptance reported; final proceedings version not established.
Cited version: arXiv 2510.02373v1 (29 September 2025).
\url{https://arxiv.org/abs/2510.02373v1}.

\leavevmode\vadjust pre{\hypertarget{ref-Zhan2026AuthMemBench}{}}%
{[}62{]} Q. Zhan, R. Zhang, S. Guo, L. Zhao, and Z. Liu. 2026. When
Memory Becomes Authority: Benchmarking Authority Collapse at the Memory
Consolidation Boundary. arXiv preprint 2608.01679v2.
\url{https://arxiv.org/abs/2608.01679v2}.

\leavevmode\vadjust pre{\hypertarget{ref-Xu2026PPMF}{}}%
{[}63{]} J. Xu, Y. Xiao, W. Shao, H. Liu, and X. Li. 2026. Memory
Provenance Laundering in LLM Agents: A Non-Amplification Firewall for
Persistent Memory. arXiv preprint 2607.29167v1.
\url{https://arxiv.org/abs/2607.29167v1}.

\leavevmode\vadjust pre{\hypertarget{ref-Lam2026SSGM}{}}%
{[}64{]} C. Lam, J. Li, L. Zhang, and K. Zhao. 2026. Governing Evolving
Memory in LLM Agents: Risks, Mechanisms, and the Stability and Safety
Governed Memory (SSGM) Framework. arXiv preprint 2603.11768v2.
\url{https://arxiv.org/abs/2603.11768v2}.

\leavevmode\vadjust pre{\hypertarget{ref-Margalit2026GovernedMemory}{}}%
{[}65{]} Y. Margalit, N. Cohen-Inger, E. Avram, R. Taig, and O.
Margalit. 2026. Governed Shared Memory for Multi-Agent LLM Systems.
arXiv preprint 2606.24535v1. \url{https://arxiv.org/abs/2606.24535v1}.

\leavevmode\vadjust pre{\hypertarget{ref-Souza2025ProvAgent}{}}%
{[}66{]} R. Souza, A. Gueroudji, S. DeWitt, D. Rosendo, T. Ghosal, R.
Ross, P. Balaprakash, and R. Ferreira da Silva. 2025. PROV-AGENT:
Unified Provenance for Tracking AI Agent Interactions in Agentic
Workflows. \emph{Proceedings of the 21st IEEE International Conference
on e-Science}. Cited version: arXiv 2508.02866v3 (20 August 2025).
\url{https://arxiv.org/abs/2508.02866v3}.

\leavevmode\vadjust pre{\hypertarget{ref-Kaufman2011Leakage}{}}%
{[}67{]} S. Kaufman, S. Rosset, and C. Perlich. 2011. Leakage in Data
Mining: Formulation, Detection, and Avoidance. \emph{Proceedings of the
17th ACM SIGKDD International Conference on Knowledge Discovery and Data
Mining}, pp. 556--563. \url{https://doi.org/10.1145/2020408.2020496}.

\end{CSLReferences}

\end{document}
